\documentclass[10pt,a4paper]{article}
\usepackage{fontspec}
\usepackage{polyglossia}
\setmainlanguage{russian}
\setotherlanguage{english}
\setmainfont{DejaVu Serif}
\setsansfont{DejaVu Sans}
\setmonofont{DejaVu Sans Mono}
\usepackage{amsmath,amssymb}
\usepackage[margin=2.2cm]{geometry}
\usepackage{titlesec}
\usepackage{enumitem}
\usepackage{multirow}
\usepackage{xcolor}
\usepackage{tcolorbox}
\tcbuselibrary{skins,breakable}
\usepackage{hyperref}
\hypersetup{colorlinks=true, linkcolor=blue!50!black, citecolor=blue!50!black, urlcolor=blue!50!black}
\usepackage{fancyhdr}
\usepackage{booktabs}
\usepackage{array}
\usepackage{tocloft}

\renewcommand{\cftsecleader}{\cftdotfill{\cftdotsep}}
\renewcommand{\cftsubsecleader}{\cftdotfill{\cftdotsep}}
\renewcommand{\cftsubsubsecleader}{\cftdotfill{\cftdotsep}}

\renewcommand{\cftsecpagefont}{\normalfont}
\renewcommand{\cftsubsecpagefont}{\normalfont}
\renewcommand{\cftsubsubsecpagefont}{\normalfont}

\setlength{\cftsecnumwidth}{2.5em}
\setlength{\cftsubsecnumwidth}{3.2em}
\setlength{\cftsubsubsecnumwidth}{4.2em}

\setlength{\cftsecindent}{0em}
\setlength{\cftsubsecindent}{1.5em}
\setlength{\cftsubsubsecindent}{3em}

\definecolor{mainblue}{RGB}{30,60,110}
\definecolor{boxbg}{RGB}{237,242,250}
\definecolor{boxframe}{RGB}{140,170,210}
\definecolor{intuitbg}{RGB}{255,247,230}
\definecolor{intuitframe}{RGB}{220,180,90}

\titleformat{\section}
  {\normalfont\Large\bfseries\color{mainblue}\raggedright}
  {\thesection}{1em}{}[{\color{mainblue}\titlerule[1pt]}]
\titlespacing*{\section}{0pt}{3.5ex plus 1ex minus .2ex}{2.3ex plus .2ex}

\titleformat{\subsection}
  {\normalfont\large\bfseries\color{mainblue}}
  {\thesubsection}{1em}{}
\titlespacing*{\subsection}{0pt}{3ex plus 1ex minus .2ex}{1.5ex plus .2ex}

\titleformat{\subsubsection}
  {\normalfont\normalsize\bfseries\color{mainblue}}
  {\thesubsubsection}{1em}{}
\titlespacing*{\subsubsection}{0pt}{2.5ex plus 1ex minus .2ex}{1ex plus .2ex}

\newtcolorbox{formulabox}{colback=boxbg,colframe=boxframe,boxrule=0.6pt,arc=2mm,left=6pt,right=6pt,top=4pt,bottom=4pt,breakable}
\newtcolorbox{intuition}{colback=intuitbg,colframe=intuitframe,boxrule=0.6pt,arc=2mm,left=6pt,right=6pt,top=4pt,bottom=4pt,breakable,title={\textbf{Интуиция}},fonttitle=\bfseries\color{black},coltitle=black}

\pagestyle{fancy}
\fancyhf{}
\fancyhead[L]{\small\textit{Макроэкономика финансовых рынков. Лекции}}
\fancyhead[R]{\small\textit{ИП ФЭН, 2026-2027}}
\fancyfoot[C]{\small\thepage}     
\renewcommand{\headrulewidth}{0.3pt}
\sloppy

\begin{document}
\begin{titlepage}
    \thispagestyle{empty}
    \centering

    \textcolor{mainblue}{\rule{\textwidth}{1.2pt}}

    \vspace{1.2cm}

    {\color{mainblue}\Huge\bfseries Макроэкономика\\[4pt] финансовых рынков\par}

    \vspace{0.8cm}

    {\color{mainblue}\Large Модуль 1\\[2pt] Динамические модели общего равновесия\par}

    \vspace{0.6cm}

    \textcolor{mainblue}{\rule{\textwidth}{0.4pt}}

    \vfill

    \begin{minipage}{0.85\textwidth}
        \small\itshape
        \textbf{Чо это такое?} Данный конспект подготовлен студенткой ОП "Экономика" Александрой Подлеснюк. Он содержит выкладки материала из оригиналов лекций курса по макро-2

        \medskip
        Автор честно пытался как следует расписать корректно любимую всеми экономическую интуицию, но всё же он хейтер макры, поэтому к жёлтым рамочкам лучше относиться скептически
    \end{minipage}

    \vspace{1cm}

    \textcolor{mainblue}{\rule{\textwidth}{1.2pt}}
\end{titlepage}

\newpage
\tableofcontents
\newpage

\section{Введение}
\subsection{Мотивация курса}
 
Зачем нужны строгие теории потребления и инвестиций?
\begin{itemize}
\item Потребление составляет около 2/3 ВВП;
\item Инвестиции — наиболее волатильный компонент ВВП;
\item \dots следовательно, потребление и инвестиции являются ключом к пониманию деловых циклов;
\item Норма сбережений не является экзогенной величиной — она определяется оптимизирующим поведением экономических агентов.
\end{itemize}
 
Почему мировые финансовые рынки важны даже для тех, кто не торгует активами?
\begin{itemize}
\item Почти все рецессии начинаются с финансовых обвалов;
\item Финансы — главный объект глобализации.
\end{itemize}
 
Каковы связи между финансовым сектором и реальной экономикой?
\begin{itemize}
\item Цены активов реагируют на новости о потреблении, занятости и т.д.;
\item Потребление и инвестиции, в свою очередь, зависят от цен активов.
\end{itemize}
 
\subsection{Эмпирические факты: динамика ключевых переменных}
 
\subsubsection{Динамика потребления}
 
Реальное потребление в США демонстрирует устойчивый растущий тренд и заметную сглаженность траектории. Возникает естественный вопрос: \textbf{почему динамика потребления настолько сглажена (smooth)?} В то же время в кризисные периоды (например, в 2008--2009 гг.) потребление демонстрирует резкое и отчётливое падение — как показывает динамика реального потребления в США за 2006--2010 гг., так и аналогичные ряды по России. Общий вывод из данных: \textbf{потребление резко падает во время кризисов}, что расходится с представлением о его исключительной гладкости в обычные периоды.
 
\subsubsection{Потребительская уверенность}
 
Индекс потребительских настроений Мичиганского университета (University of Michigan Consumer Sentiment Index) вместе с темпами роста ВВП и датировкой рецессий показывает, что резкие провалы потребительской уверенности систематически совпадают с рецессиями (например, начало 1980-х, 1990--1991, 2001, 2008--2009, 2020).
 
\subsubsection{Динамика инвестиций}
 
Инвестиции в США значительно более волатильны, чем потребление. \textbf{Почему инвестиции так волатильны?} — это второй ключевой эмпирический вопрос курса. Сопоставление темпов роста потребления (сглаженная кривая) и инвестиций (более "рваная", размашистая кривая) наглядно иллюстрирует этот контраст в волатильности.
 
\subsubsection{Фондовый рынок}
 
\textbf{Внутридневная динамика.} На примере акций Cisco Systems (5 апреля, внутридневные котировки) видно, что цена акции колеблется в течение одного торгового дня — возникает вопрос: \textbf{почему цены акций так волатильны?}
 
\textbf{Долгосрочная динамика.} Индекс Dow Jones с 1901 г. демонстрирует выраженный растущий тренд, но с периодическими резкими падениями — наиболее известное из них произошло во время Великой депрессии (индекс упал с пика около 380 в 1929 г. до примерно 41 в 1932 г.). Аналогичная картина волатильности с резкими коррекциями наблюдается по индексу NASDAQ Composite (пузырь доткомов конца 1990-х -- начала 2000-х: рост с уровня около 700 в 1993 г. до пика около 5000 в 2000 г. и последующее падение примерно до 1100 в 2002 г.) и по российскому индексу РТС.
 
\textbf{Дивиденды и цены акций.} Сопоставление индекса S\&P 500 и скользящей за 12 месяцев прибыли на акцию (12M OEPS) ставит вопрос: \textbf{объясняют ли дивиденды (прибыли) динамику цен акций?} — визуально цены акций демонстрируют существенно большую волатильность, чем фундаментальные показатели прибыли.
 
\textbf{P/E и процентные ставки.} Долгосрочный график коэффициента цена/прибыль (Price-Earnings Ratio, CAPE) в сопоставлении с долгосрочными процентными ставками (1880--2020-е) показывает несколько выраженных пиков переоценённости рынка — 1901, 1929, 1966, 1981, 2000 гг., что часто ассоциируется с последующими коррекциями.
 
\subsubsection{Цены на жильё}
 
Сравнение японского и американского жилищных пузырей (US Housing Bubble Timeline vs Japanese Housing Bubble Timeline) показывает похожую динамику "пузырь--схлопывание": как в Японии конца 1980-х, так и в США середины 2000-х наблюдался стремительный рост цен на жильё с последующим резким падением. \textbf{"Потерянное десятилетие" в Японии началось именно со схлопывания пузыря на рынке жилья} — иллюстрация: цена жилой земли в районе Минато-ку (Токио) выросла с уровня около 300 тыс. иен/кв.м в середине 1970-х до пика свыше 6,5 млн иен/кв.м к 1987 г., после чего обвалилась до уровня около 1 млн иен/кв.м к концу 1990-х -- 2000-м гг.
 
Цены на жильё в Москве (индекс стоимости жилья, \$/кв.м, по данным www.irn.ru) также демонстрируют выраженный пузырь середины 2000-х (пик около 6100 \$/кв.м к 2008 г.) с последующей коррекцией и стабилизацией на уровне около 2500--3000 \$/кв.м.
 
Долгосрочный американский график цен на жильё (Home Prices) в сопоставлении со строительными издержками (Building Costs), численностью населения (Population) и процентными ставками (Interest Rates) за 1880--2020 гг. показывает, что цены на жильё в период 2000-х гг. оторвались от динамики строительных издержек и населения существенно сильнее, чем в предыдущие десятилетия, — характерный признак пузыря.
 
\subsection{Ключевые вопросы курса}
 
На основе рассмотренных фактов формулируется список вопросов, на которые курс призван дать теоретические ответы:
 
\begin{itemize}
\item Что определяет потребление, сбережения и инвестиции?
\item Почему потребление сглажено, а инвестиции волатильны?
\item Почему финансовые обвалы вызывают резкое падение (реальных) частных расходов?
\item Должно ли государство наращивать заимствования и расходы, когда частный сектор сберегает слишком много?
\end{itemize}
 
\textit{Финансовая нестабильность:}
\begin{itemize}
\item Почему цены активов и валютные курсы настолько волатильны?
\item Что вызывает пузыри активов? Что вызывает обвалы (meltdowns)?
\item Что вызывает масштабные притоки и оттоки капитала?
\item Существует ли (монетарное/фискальное/макропруденциальное) политическое "лекарство" от финансовой нестабильности?
\end{itemize}
 
\newpage
\section{Лекции 1-2. Межвременной выбор}
 
\subsection{Мотивация}
 
Необходимы микрооснования для решений о потреблении и сбережениях: почти все современные макроэкономические модели строятся именно на микрооснованиях. Стандартная микроэкономика даёт статическую картину выбора, тогда как сбережения по своей природе являются динамическим феноменом.
 
Ключевой вопрос: какие факторы, помимо текущего располагаемого дохода, определяют потребление и сбережения (а значит, и потребительскую уверенность)? В частности — оказывает ли фискальное стимулирование какой-либо эффект на потребление во время рецессии?
 
\subsection{Суть современных теорий потребления}
 
\begin{itemize}
\item Потребительский выбор является \textbf{межвременным}: он учитывает текущее и будущие состояния экономики.
\item Потребители формируют \textbf{рациональные ожидания} относительно будущих состояний.
\item Отправная точка — упрощающее предположение о \textbf{совершенном знании}: потребители знают весь свой будущий доход (горизонт жизни, процентные ставки и т.д.).
\item Потребление — \textbf{forward-looking} (устремлённая в будущее) переменная.
\end{itemize}
 
Динамический выбор строится на двух компонентах:
\begin{itemize}
\item пожизненная полезность от потребления в разные периоды;
\item динамическое (межвременное) бюджетное ограничение.
\end{itemize}
 
\subsection{Пожизненная полезность}
 
Рассмотрим потребителя с горизонтом жизни $T$ периодов, максимизирующего
\begin{equation}
\sum_{t=0}^{T-1} \frac{u(C_t)}{(1+\rho)^t} \longrightarrow \max_{C_t} \tag{2.1}
\end{equation}
 
Потребление в периоде $t$ приносит мгновенную полезность $u(C_t)$ со стандартными неоклассическими свойствами: $u'(C) > 0$, $u''(C) < 0$ (предельная полезность положительна и убывает). Мгновенные предпочтения не меняются во времени.
 
Потребитель нетерпелив и дисконтирует будущий поток полезности по ставке $\rho$ (субъективная норма дисконтирования). Альтернативное обозначение — субъективный дисконт-фактор $\beta = 1/(1+\rho)$.
 
\begin{intuition}
Чем выше $\rho$, тем сильнее потребитель обесценивает будущую полезность относительно текущей, то есть тем более он "нетерпелив". Предположения $u'>0$ и $u''<0$ гарантируют, что дополнительная единица потребления всегда желательна, но её ценность убывает по мере роста уровня потребления — это свойство лежит в основе желания сглаживать потребление во времени.
\end{intuition}
 
\subsubsection*{Отступление: временн\'ая несогласованность предпочтений (time-inconsistent preferences)}
 
Субъективное дисконтирование будущей полезности — поведенческий феномен. Хотя геометрическое дисконтирование удобно для моделирования, оно может быть нереалистичным. В действительности люди сильнее обесценивают слишком отдалённое будущее. Например, можно предположить, что потребитель различает только "сегодня" и "всё остальное будущее":
\[
U = u(C_0) + \frac{1}{1+\rho}\sum_{t=1}^{T-1} u(C_t)
\]
 
Более реалистичным считается \textbf{гиперболическое} субъективное дисконтирование, однако оно существенно усложняет алгебру модели:
\[
\rho(\tau) = \frac{1}{1+\alpha\tau}
\]
где $\tau$ — горизонт (расстояние во времени), а $\alpha$ — параметр, определяющий скорость затухания "веса" дисконтирования по мере удаления в будущее.
 
\subsection{Динамическое (потоковое) бюджетное ограничение}
 
Введём обозначения:
\begin{itemize}
\item $Y_t^d$ — располагаемый доход в периоде $t$: доход за вычетом чистых паушальных налогов, $Y_t^d = Y_t - T_t$;
\item $S_t = Y_t^d - C_t$ — сбережения в периоде $t$ (это \textbf{потоковая} переменная!);
\item $A_t$ — активы \textit{на начало} периода $t$ (положительны, если это активы, отрицательны, если это долг) — это \textbf{запасовая} переменная!
\item $r$ — процентная ставка, предполагается постоянной для упрощения обозначений (это не принципиально); первоначально предполагается отсутствие ограничений ликвидности.
\end{itemize}
 
\begin{align}
A_{t+1} &= (A_t + S_t)(1+r) \notag \\
&= (A_t + Y_t^d - C_t)(1+r) \tag{2.2}
\end{align}
 
\begin{intuition}
Уравнение (2.2) — это стандартное правило накопления активов: то, чем потребитель располагает в начале следующего периода, равно тому, что было отложено (или занято) в текущем периоде, с учётом начисленных на всю сумму процентов. Сбережение $S_t$ — это \emph{поток} за период, тогда как $A_t$ — \emph{запас} на определённый момент времени; важно не путать эти две категории переменных.
\end{intuition}
 
\subsection{Задача оптимизации: составление Лагранжиана}
 
Потребитель максимизирует (2.1) при ограничении (2.2). Составим функцию Лагранжа:
\[
\Lambda = \sum_{t=0}^{T-1} \frac{u(C_t)}{(1+\rho)^t} + \sum_{t=0}^{T-1} \lambda_t \Big[A_{t+1} - (A_t + Y_t^d - C_t)(1+r)\Big]
\]
 
Условия первого порядка:
\begin{align}
\frac{\partial \Lambda}{\partial C_t} &= \frac{u'(C_t)}{(1+\rho)^t} + \lambda_t(1+r) = 0 \notag\\
\frac{\partial \Lambda}{\partial C_{t+1}} &= \frac{u'(C_{t+1})}{(1+\rho)^{t+1}} + \lambda_{t+1}(1+r) = 0 \notag\\
\frac{\partial \Lambda}{\partial A_{t+1}} &= \lambda_t - \lambda_{t+1}(1+r) = 0 \notag
\end{align}
 
Из этих трёх условий, исключая множители Лагранжа, получаем ключевое соотношение:
\begin{equation}
\frac{u'(C_t)}{u'(C_{t+1})} = \frac{1+r}{1+\rho} \tag{2.3}
\end{equation}
 
\subsection{Уравнение Эйлера (правило Рамсея--Кейнса)}
 
Уравнение (2.3) называется уравнением Эйлера, или правилом Рамсея--Кейнса. Перепишем его в форме, явно указывающей на интерпретацию через относительные цены:
\[
\frac{u'(C_t)}{u'(C_{t+1})(1+\rho)^{-1}} = \frac{1}{(1+r)^{-1}}
\]
 
Нормализуем цену потребления в периоде $t$ к единице. Тогда $(1+r)^{-1}$ — это относительная цена потребления в периоде $t+1$.
 
\begin{intuition}
Предельная норма замещения (MRS) между потреблением в периодах $t$ и $t+1$ должна быть равна соотношению их цен. Это в точности стандартное условие оптимальности потребительского выбора из микроэкономики, только применённое не к двум разным благам, а к потреблению одного и того же блага в разные моменты времени. Левая часть (2.3) — это MRS между $C_t$ и $C_{t+1}$ с учётом субъективного дисконтирования; правая часть — рыночная относительная цена "завтрашнего" потребления в терминах "сегодняшнего".
\end{intuition}
 
\subsection{Эффект межвременного замещения}
 
Из (2.3):
\[
r \uparrow \;\Rightarrow\; \frac{u'(C_t)}{u'(C_{t+1})} \uparrow \;\Rightarrow\; C_t \downarrow,\; S_t \uparrow,\; C_{t+1}\uparrow
\]
 
\begin{intuition}
Рост процентной ставки снижает относительную цену потребления в периоде $t+1$ и увеличивает предельную норму замещения между потреблением в периодах $t$ и $t+1$. Поскольку предельная полезность убывает по потреблению, это означает снижение потребления в периоде $t$, рост сбережений в периоде $t$ и рост потребления в периоде $t+1$. Иными словами, потребитель замещает "дорогое" сегодняшнее потребление "подешевевшим" относительно него завтрашним.
\end{intuition}
 
\subsection{Эффект дохода}
 
Предположим, что $S_t > 0$. Тогда
\[
r\uparrow \;\Rightarrow\; (1+r)S_t \uparrow \;\Rightarrow\; C_t \uparrow,\; C_{t+1}\uparrow,\; S_t\downarrow
\]
 
\begin{intuition}
Рост процентной ставки увеличивает доход от капитала в периоде $t+1$ и, следовательно, увеличивает суммарный (пожизненный) доход потребителя. Обратите внимание: $Y_t$ — это поток трудового дохода (или дохода от "надела"), а не совокупный доход. Рост процентного дохода в будущем позволяет увеличить потребление \emph{в обоих} периодах — это классический эффект дохода. Полезно самостоятельно рассмотреть также случаи $S_t < 0$ (должник — рост $r$ увеличивает бремя обслуживания долга и действует в обратную сторону) и $S_t = 0$ (эффект дохода отсутствует).
\end{intuition}
 
\subsection{Свойства сбережений: эластичность межвременного замещения}
 
Эффект межвременного замещения и эффект дохода действуют в противоположных направлениях в отношении сбережений (и потребления) в периоде $t$. Итоговый эффект зависит от \textbf{межвременной (мгновенной) эластичности замещения}:
\[
\sigma = -\frac{u'(c)}{u''(c)\,c} = \lim_{\tau\to t}\sigma(c_t, c_\tau), \qquad
\sigma(c_t,c_\tau) = -\frac{d\ln(c_\tau/c_t)}{d\ln MRS(c_t,c_\tau)}
\]
 
\[
\frac{\partial S_t}{\partial r}
\begin{cases}
>0, & \text{если } \sigma > 1 \quad \text{(этот случай будет предполагаться далее)}\\
=0, & \text{если } \sigma = 1\\
<0, & \text{если } \sigma < 1
\end{cases}
\]
 
\begin{intuition}
Параметр $\sigma$ измеряет, насколько сильно потребитель готов менять соотношение потребления в разные периоды в ответ на изменение относительной цены (то есть процентной ставки). При высокой эластичности замещения ($\sigma>1$) эффект замещения преобладает над эффектом дохода, и рост $r$ увеличивает сбережения. При низкой эластичности ($\sigma<1$) доминирует эффект дохода, и рост $r$, наоборот, сокращает сбережения (потребитель может "позволить себе" меньше сберегать, чтобы достичь того же уровня будущего потребления). Далее в курсе принимается допущение $\sigma>1$ как базовый случай.
\end{intuition}
 
\subsection{Динамика потребления во времени}
 
Из уравнения Эйлера (2.3) следует:
\begin{itemize}
\item Если $r > \rho$ — потребление растёт на протяжении жизни;
\item Если $r < \rho$ — потребление снижается на протяжении жизни;
\item Если $r = \rho$ — оптимальным выбором является постоянное потребление.
\end{itemize}
 
\begin{intuition}
\textbf{Важно:} эти динамические свойства \emph{не зависят} от того, растёт ли, падает ли или постоянен доход во времени. Профиль потребления определяется исключительно соотношением рыночной доходности $r$ и субъективной нормы дисконтирования $\rho$, тогда как \emph{уровень} потребления в каждом периоде определяется бюджетным ограничением (совокупным богатством). Это разделение — ключевая особенность модели межвременного выбора с рациональными ожиданиями.
\end{intuition}
 
\subsection{От динамического к межвременному бюджетному ограничению}
 
Из (2.2) выразим $A_t$ и подставим рекурсивно вперёд:
\[
A_t = \frac{A_{t+1}}{(1+r)} - S_t, \qquad A_{t+1} = \frac{A_{t+2}}{(1+r)} - S_{t+1}
\]
\[
A_t = -S_t - \frac{S_{t+1}}{(1+r)} + \frac{A_{t+2}}{(1+r)^2} = -S_t - \frac{S_{t+1}}{(1+r)} - \frac{S_{t+2}}{(1+r)^2} + \frac{A_{t+3}}{(1+r)^3} = \dots
\]
 
Продолжая итерации вплоть до конца горизонта жизни $T$:
\begin{equation}
A_t = -\sum_{\tau=0}^{T-t-1}\frac{S_{t+\tau}}{(1+r)^\tau} + \frac{A_T}{(1+r)^{T-t}} \tag{в исходной форме}
\end{equation}
 
\subsection{Условие отсутствия "схемы Понци" (No-Ponzi game / трансверсальность)}
 
\begin{equation}
\frac{A_T}{(1+r)^{T-t}} = 0 \tag{2.4}
\end{equation}
 
\begin{intuition}
Рациональный потребитель не должен оставлять после смерти никаких активов (при отсутствии мотива наследования — bequest motive). Симметрично, рациональные кредиторы не должны предоставлять кредит потребителю, который не собирается его возвращать. Условие (2.4) — это ограничение, накладываемое на \emph{будущий} выбор потребителя (в отличие от уравнения накопления активов (2.2), которое смотрит "назад", в прошлое).
\end{intuition}
 
Подставляя условие (2.4) в итерированное уравнение накопления активов, получаем:
\begin{equation}
A_t = -\sum_{\tau=0}^{T-t-1} \frac{S_{t+\tau}}{(1+r)^\tau} \tag{2.5}
\end{equation}
 
\subsection{Интерпретация межвременного бюджетного ограничения}
 
\textbf{Случай $A_t > 0$:}
\[
A_t > 0 \;\Rightarrow\; \sum_{\tau=0}^{T-t-1} \frac{S_{t+\tau}}{(1+r)^\tau} < 0
\]
Владение активами сегодня допускает отрицательную приведённую стоимость будущих сбережений — т.е. потребитель может тратить больше своего (трудового) дохода, потребляя как из дохода, так и из накопленного богатства.
 
\textbf{Случай $A_t < 0$:}
\[
A_t < 0 \;\Rightarrow\; \sum_{\tau=0}^{T-t-1} \frac{S_{t+\tau}}{(1+r)^\tau} > 0
\]
Задолженность сегодня должна быть обеспечена положительной приведённой стоимостью будущих сбережений: потребитель должен тратить меньше своего дохода, используя доход для потребления и погашения долга.
 
\begin{intuition}
Условие NPG — это ограничение на \emph{будущее} поведение потребителя, тогда как динамика $A_t$ является \emph{backward-looking} (определяется прошлым). Именно сочетание этих двух свойств не позволяет потребителю бесконечно наращивать долг, "перезанимая" для обслуживания старого долга (собственно, "схема Понци").
\end{intuition}
 
\subsection{Межвременное бюджетное ограничение}
 
Подставим определение $S_{t+\tau} = Y_{t+\tau}^d - C_{t+\tau}$ в (2.5):
\begin{align}
A_t &= -\sum_{\tau=0}^{T-t-1}\frac{S_{t+\tau}}{(1+r)^\tau} = -\sum_{\tau=0}^{T-t-1}\frac{Y_{t+\tau}^d - C_{t+\tau}}{(1+r)^\tau}\notag\\
\sum_{\tau=0}^{T-t-1}\frac{C_{t+\tau}}{(1+r)^\tau} &= A_t + \sum_{\tau=0}^{T-t-1}\frac{Y_{t+\tau}^d}{(1+r)^\tau} \tag{2.6}
\end{align}
 
\begin{intuition}
Уравнение (2.6) — центральный результат этого раздела: в каждом периоде приведённая (дисконтированная) стоимость всех будущих расходов на потребление должна быть равна сумме текущего богатства $A_t$ и приведённой стоимости всех будущих потоков располагаемого дохода. Именно эта величина — совокупное "пожизненное богатство" — и определяет уровень потребления, тогда как уравнение Эйлера определяет лишь его динамику (профиль во времени).
\end{intuition}
 
\subsubsection*{Отступление: альтернативная спецификация задачи потребителя}
 
Можно показать, что альтернативная постановка задачи потребителя даёт то же самое условие первого порядка / уравнение Эйлера (2.3):
\[
\sum_{t=0}^{T-1}\frac{u(C_t)}{(1+\rho)^t} \to \max_{C_t}\quad\text{при}\quad \sum_{t=0}^{T-1}\frac{C_t}{(1+r)^t} = A_0 + \sum_{t=0}^{T-1}\frac{Y_t^d}{(1+r)^t} \tag{2.6'}
\]
Эта формулировка требует всего одного множителя Лагранжа (в отличие от последовательности $\lambda_t$ в исходной постановке), поскольку ограничение сразу задано в межвременной, а не в потоковой форме.
 
\subsection{Условие NPG и бюджетное ограничение при бесконечном горизонте}
 
При $T\to\infty$ условие отсутствия "схемы Понци" принимает вид:
\begin{equation}
\lim_{T\to\infty}\frac{A_T}{(1+r)^{T-t}} = 0 \tag{2.7}
\end{equation}
 
\begin{intuition}
Приведённая стоимость активов/долга на бесконечном горизонте равна нулю; активы (или долг) не должны расти темпом, превышающим рыночную процентную ставку.
\end{intuition}
 
Соответственно, потоковое бюджетное ограничение при бесконечном горизонте:
\begin{equation}
A_t = -\sum_{\tau=0}^{\infty}\frac{S_{t+\tau}}{(1+r)^\tau} \tag{2.8}
\end{equation}
 
и межвременное бюджетное ограничение:
\begin{equation}
\sum_{\tau=0}^{\infty}\frac{C_{t+\tau}}{(1+r)^\tau} = A_t + \sum_{\tau=0}^{\infty}\frac{Y_{t+\tau}^d}{(1+r)^\tau} \tag{2.9}
\end{equation}
 
\subsection{Кейс: падение потребительской уверенности в 2008 г.}
 
Индекс потребительской уверенности в США (United States Consumer Confidence) демонстрирует резкое падение в 2008 г. — с уровня выше 100--140 в конце 1990-х -- середине 2000-х до примерно 25--30 в начале 2009 г. С точки зрения межвременного бюджетного ограничения (2.6)/(2.9), резкое падение желаемого текущего потребления должно быть связано либо с падением текущего богатства $A_t$, либо с пересмотром вниз ожидаемых будущих потоков располагаемого дохода $Y_{t+\tau}^d$, либо с ростом ставки дисконтирования будущих доходов.
 
Возможные объяснения резкого падения потребительской уверенности в 2008 г., предлагаемые для обсуждения:
\begin{itemize}
\item Лопнувший пузырь на рынке жилья (падение стоимости активов домохозяйств, $A_t\downarrow$);
\item Обвал на фондовом рынке (аналогично, падение $A_t$);
\item Рост ставок по кредитам (изменение условий дисконтирования / доступности заимствований);
\item Страх потери работы (или сокращения зарплаты) — то есть пересмотр вниз ожидаемых будущих $Y_{t+\tau}^d$.
\end{itemize}
 
Отдельно ставится ключевой вопрос: \textbf{почему фискальное стимулирование не помогло?} Этот вопрос иллюстрируется карикатурой из журнала The Economist (1 мая 2008 г.): Федеральная резервная система последовательно снижает процентную ставку ("Time to drop the interest rate again"), однако экономика ("Economy"661), которую тянет за собой фигура "US Consumers", обременённая долгом, не демонстрирует движения ("Any movement?"). Иллюстрация подчёркивает идею, что одного лишь монетарного/фискального стимула может быть недостаточно, если пересмотрены вниз ожидания домохозяйств относительно будущего дохода или богатства — именно на этот вопрос отвечает анализ через межвременное бюджетное ограничение и рикардианскую эквивалентность (см. далее).
 
\subsection{Динамическое бюджетное ограничение государства}
 
Введём обозначения, аналогичные обозначениям для домохозяйства:
\begin{itemize}
\item $G_t$ — государственные закупки в периоде $t$;
\item $T_t$ — чистые паушальные налоги в периоде $t$;
\item $b_t$ — государственный долг на начало периода $t$ (положителен, если это долг, отрицателен, если это государственные активы);
\item $d_t = G_t - T_t$ — первичный дефицит бюджета в периоде $t$ ("первичный" означает без учёта обслуживания долга);
\item $r$ — процентная ставка, как и прежде постоянная и совпадающая со ставкой для частных агентов.
\end{itemize}
 
\begin{align}
b_{t+1} &= (b_t + d_t)(1+r)\notag\\
&= (b_t + G_t - T_t)(1+r) \tag{2.10}
\end{align}
 
Итерируя (2.10) вперёд так же, как и в случае домохозяйства:
\[
b_t = \frac{b_{t+1}}{(1+r)} - d_t,\qquad b_{t+1} = \frac{b_{t+2}}{(1+r)} - d_{t+1}
\]
\[
b_t = -d_t - \frac{d_{t+1}}{(1+r)} + \frac{b_{t+2}}{(1+r)^2} = -d_t - \frac{d_{t+1}}{(1+r)} - \frac{d_{t+2}}{(1+r)^2} + \frac{b_{t+3}}{(1+r)^3} = \dots
\]
\[
b_t = -\sum_{\tau=0}^{T-t-1}\frac{d_{t+\tau}}{(1+r)^\tau} + \frac{b_T}{(1+r)^{T-t}}
\]
 
\subsubsection{NPG-условие для государственного долга}
 
\begin{equation}
\lim_{T\to\infty}\frac{b_T}{(1+r)^{T-t}} = 0 \tag{2.11}
\end{equation}
 
\begin{intuition}
Приведённая стоимость государственного долга на бесконечном горизонте равна нулю; государственный долг не должен расти "слишком быстро" — темпом, превышающим процентную ставку. Заметим: государство, естественно, имеет бесконечный горизонт планирования (в отличие от домохозяйства с конечным $T$), поэтому условие (2.11) формулируется сразу для $T\to\infty$.
\end{intuition}
 
\subsubsection{От динамического к межвременному бюджетному ограничению государства}
 
Аналогично случаю домохозяйства, из (2.10) и (2.11) получаем:
\begin{equation}
b_t = -\sum_{\tau=0}^{\infty}\frac{d_{t+\tau}}{(1+r)^\tau} \tag{2.12}
\end{equation}
 
Подставляя $d_{t+\tau} = G_{t+\tau} - T_{t+\tau}$:
\begin{equation}
\sum_{\tau=0}^{\infty}\frac{G_{t+\tau}}{(1+r)^\tau} + b_t = \sum_{\tau=0}^{\infty}\frac{T_{t+\tau}}{(1+r)^\tau} \tag{2.13}
\end{equation}
 
\begin{intuition}
В каждом периоде сумма приведённой стоимости будущих государственных расходов и текущего государственного долга должна быть равна приведённой стоимости будущих налоговых поступлений:
\[
b_t > 0 \;\Rightarrow\; \sum_{\tau=0}^{\infty}\frac{d_{t+\tau}}{(1+r)^\tau} < 0
\]
Государственный долг должен быть обеспечен соответствующей приведённой стоимостью будущих бюджетных профицитов — это \textbf{принцип устойчивой (sustainable) фискальной политики}. Условие NPG для государственного долга — это ограничение на будущую фискальную политику, тогда как динамика долга является backward-looking (определяется прошлым), в полной аналогии со случаем домохозяйства.
\end{intuition}
 
\subsubsection{Долговая нагрузка: эмпирическая иллюстрация}
 
Материалы лекции иллюстрируют динамику отношения государственного долга к ВВП в США (заметный рост, особенно после кризиса 2008 г. и в период пандемии), а также глобальную картину отношения валового государственного долга к ВВП по состоянию на 2018 г. (данные МВФ, World Economic Outlook, октябрь 2017 г.): наиболее высокая долговая нагрузка (100\% ВВП и более) характерна для США, Японии и ряда стран Южной Европы, тогда как многие развивающиеся страны, а также Россия, имеют относительно низкий уровень государственного долга (менее 25\% ВВП).
 
\subsection{Объединение межвременных бюджетных ограничений: рикардианская эквивалентность}
 
Предположим, что государственные облигации являются единственным видом активов в экономике, то есть $A_t = b_t$. Сопоставим межвременное бюджетное ограничение домохозяйства (2.9) (записанное с явным учётом налогов $T_{t+\tau}$ вместо обобщённого $Y^d$):
\begin{equation}
\sum_{\tau=0}^{\infty}\frac{C_{t+\tau}}{(1+r)^\tau} = A_t + \sum_{\tau=0}^{\infty}\frac{(Y_{t+\tau}-T_{t+\tau})}{(1+r)^\tau} \tag{2.9}
\end{equation}
и межвременное бюджетное ограничение государства (2.13). Подставляя (2.13) в (2.9) (то есть заменяя $A_t = b_t$ выражением из (2.13)), приходим к:
\begin{equation}
\sum_{\tau=0}^{\infty}\frac{C_{t+\tau}}{(1+r)^\tau} = \sum_{\tau=0}^{\infty}\frac{(Y_{t+\tau}-G_{t+\tau})}{(1+r)^\tau} \tag{2.14}
\end{equation}
 
\begin{intuition}
Это и есть формальный вывод \textbf{рикардианской эквивалентности}: после объединения бюджетных ограничений домохозяйства и государства переменные $T_t$ и $b_t$ полностью исчезают из совокупного (межвременного) бюджетного ограничения — совокупное потребление зависит только от приведённой стоимости будущего дохода за вычетом государственных расходов, но \emph{не} от того, как именно эти расходы финансируются (налогами сейчас или долгом, то есть налогами позже).
\end{intuition}
 
\subsubsection{Формулировка результата рикардианской эквивалентности}
 
\textbf{Способ, которым государство финансирует свои расходы — за счёт налогов или за счёт долга — не влияет на решения о потреблении.} Имеет значение только приведённая (дисконтированная) стоимость будущих государственных расходов. При устойчивой фискальной политике снижение налогов сегодня требует повышения налогов в будущем — для сохранения адекватного обеспечения государственного долга.
 
Классическая ссылка: Robert Barro, "Are Government Bonds Net Wealth?" — ответ: \textbf{нет!}
\begin{center}
\textit{Barro R.\,J. (1974) ``Are Government Bonds Net Wealth?'' Journal of Political Economy, 82(6), pp.\,1095--1117.}
\end{center}
 
\subsubsection{Факторы, нарушающие рикардианскую эквивалентность}
 
\begin{itemize}
\item Конечные горизонты жизни и отсутствие межпоколенческого альтруизма;
\item Финансовые несовершенства / ограничения ликвидности — в частности, различие процентных ставок для частных агентов и для государства;
\item Искажающее (distortionary) налогообложение.
\end{itemize}
 
\begin{intuition}
Если горизонт жизни домохозяйства конечен, а государство фактически бессмертно (постоянно рефинансирует долг за счёт последующих поколений), то бремя будущих налогов может лечь на поколения, которые не получают выгоды от сегодняшнего снижения налогов, — тогда снижение налогов, финансируемое долгом, воспринимается домохозяйствами как чистый прирост богатства, и рикардианская эквивалентность нарушается. Аналогично, если домохозяйства сталкиваются с ограничениями по заимствованию (не могут занимать под будущий доход так же свободно, как государство), снижение налогов сегодня ослабляет эти ограничения и стимулирует потребление. Наконец, если налоги являются искажающими (например, зависят от уровня дохода или капитала), то их межвременное перераспределение влияет на реальные стимулы (к труду, инвестированию), а не только на распределение богатства во времени.
\end{intuition}
 
\subsection{Фискальные мультипликаторы}
 
\subsubsection{Оценка Ромера и Бернстайна (2009)}
 
По расчётам Romer and Bernstein (2009), динамика фискальных мультипликаторов на горизонте до 16 кварталов показывает, что мультипликатор увеличения государственных расходов (Spending Increase) стабилизируется на уровне около 1{,}57, тогда как мультипликатор снижения налогов (Tax Cut) стабилизируется на более низком уровне — около 0{,}99. Оба мультипликатора нарастают постепенно, достигая "плато" примерно через 6--8 кварталов.
 
\begin{intuition}
Более низкий мультипликатор налогового стимула по сравнению с мультипликатором прямых государственных расходов согласуется с логикой межвременного выбора: часть дополнительного располагаемого дохода от снижения налогов домохозяйства сберегают (в соответствии с уравнением Эйлера и желанием сгладить потребление), тогда как государственные закупки товаров и услуг напрямую входят в совокупный спрос.
\end{intuition}
 
\subsubsection{Эмпирические оценки мультипликатора государственных расходов}
 
\begin{center}
\begin{tabular}{p{5.2cm}p{4cm}p{2.5cm}}
\toprule
\textbf{Источник} & \textbf{Страна} & \textbf{Годовой мультипликатор}\\
\midrule
Blanchard and Perotti (2002) & США & 0{,}5--0{,}6\\
Bryant и др. (1988) & США & 0{,}6--2{,}0\\
Cogan и др. (2009) & США & 0{,}7--0{,}9\\
\multirow{2}{5.2cm}{Dalsgaard, André, Richardson (2001)} & США & 1{,}1\\
& Япония & 1{,}7\\
Elmendorf and Furman (2008) & США & 1\\
\multirow{6}{5.2cm}{HM Treasury (2003)} & Германия & 0{,}4\\
& Испания & 0{,}5\\
& Франция & 0{,}5\\
& Португалия & 0{,}7\\
& Швеция & 0{,}4\\
& Великобритания & 0{,}3\\
\multirow{2}{5.2cm}{Ilzetzki and V\'egh (2008)} & Развитые страны & 0{,}7\\
& Развивающиеся страны & 0{,}4\\
\multirow{2}{5.2cm}{IMF (2008)} & Развитые страны & 0{,}2\\
& Развивающиеся страны & 0{,}1\\
\bottomrule
\end{tabular}
\end{center}
 
\subsubsection{Эмпирические оценки налогового мультипликатора}
 
\begin{center}
\begin{tabular}{p{5.2cm}p{4cm}p{2.5cm}}
\toprule
\textbf{Источник} & \textbf{Страна} & \textbf{Годовой мультипликатор}\\
\midrule
Blanchard and Perotti (2002) & США & 0{,}7--1{,}1\\
Cogan и др. (2009) & США & 0{,}7--0{,}9\\
\multirow{6}{5.2cm}{HM Treasury (2003)} & Германия & 0{,}2\\
& Испания & 0{,}1\\
& Франция & 0{,}1\\
& Португалия & 0--0{,}1\\
& Швеция & 0{,}3\\
& Великобритания & 0{,}2\\
\multirow{2}{5.2cm}{IMF (2008)} & Развитые страны & 0\\
& Развивающиеся страны & 0{,}1\\
Romer and Romer (2008) & США & 1{,}2\\
\bottomrule
\end{tabular}
\end{center}
 
\begin{intuition}
Разброс эмпирических оценок мультипликаторов между исследованиями и странами весьма велик. Тем не менее общая закономерность устойчива: мультипликатор государственных расходов в подавляющем большинстве оценок выше (или, по крайней мере, не ниже) мультипликатора снижения налогов, что согласуется с теоретическим предсказанием о частичном сбережении дополнительного располагаемого дохода домохозяйствами.
\end{intuition}
 
\subsection{Кривая Лаффера для налога на труд}
 
По расчётам Trabandt and Uhlig (2012), для стационарных калибровок моделей США и "ЕС-14" построена кривая зависимости стационарных налоговых поступлений (Steady State Tax Revenues, среднее по выборке нормировано на 100) от ставки налога на труд (Steady State Labor Tax). Обе кривые имеют классическую форму перевёрнутой параболы (кривая Лаффера):
\begin{itemize}
\item Для США максимум налоговых поступлений достигается при ставке налога на труд около 0{,}62 (пиковое значение поступлений — около 137{,}5 от среднего);
\item Для "ЕС-14" максимум — также около ставки 0{,}61--0{,}62, но при более низком пиковом уровне поступлений (около 112 от среднего);
\item Фактическое положение США на кривой (отмечено звёздочкой) соответствует ставке около 0{,}22 — то есть значительно левее пика, в зоне, где рост ставки увеличивает поступления;
\item Фактическое положение "ЕС-14" — ставка около 0{,}34, также левее соответствующего пика.
\end{itemize}
 
\begin{intuition}
Кривая Лаффера иллюстрирует фундаментальный компромисс налоговой политики: при низких ставках рост ставки налога увеличивает поступления почти пропорционально (эффект искажения стимулов к труду ещё невелик), однако по мере роста ставки предложение труда сокращается всё сильнее, и в конце концов дальнейшее повышение ставки \emph{снижает} налоговые поступления. Согласно расчётам Trabandt and Uhlig, ни США, ни "ЕС-14" в среднем не находятся на "неверной" (нисходящей) стороне кривой Лаффера — то есть, по крайней мере в рамках калибровки модели, дальнейшее повышение ставок налога на труд в среднем могло бы увеличить, а не уменьшить, налоговые поступления, хотя разрыв до пика у стран ЕС-14 меньше, чем у США.
\end{intuition}
 
\subsection{Краткое резюме ключевых результатов}
 
\begin{itemize}
\item \textbf{Эмпирическая мотивация}: потребление — около 2/3 ВВП и сглажено во времени, но резко падает в кризисы; инвестиции и цены активов (акции, жильё) крайне волатильны и подвержены пузырям и обвалам.
\item \textbf{Пожизненная полезность и уравнение Эйлера}: оптимальный межвременной выбор потребителя описывается условием $\dfrac{u'(C_t)}{u'(C_{t+1})} = \dfrac{1+r}{1+\rho}$; динамика потребления определяется соотношением $r$ и $\rho$ и не зависит от профиля дохода во времени.
\item \textbf{Эффекты процентной ставки}: рост $r$ порождает разнонаправленные эффект замещения (в пользу сбережений) и эффект дохода (в пользу текущего потребления, если сбережения положительны); итог зависит от эластичности межвременного замещения $\sigma$.
\item \textbf{Межвременное бюджетное ограничение}: приведённая стоимость всех будущих расходов на потребление равна текущему богатству плюс приведённая стоимость будущих доходов; условие NPG (отсутствие "схемы Понци") исключает бесконечное наращивание долга.
\item \textbf{Государственное бюджетное ограничение} строится симметрично бюджетному ограничению домохозяйства: долг должен быть обеспечен приведённой стоимостью будущих бюджетных профицитов (принцип устойчивой фискальной политики).
\item \textbf{Рикардианская эквивалентность}: при объединении бюджетных ограничений домохозяйства и государства способ финансирования госрасходов (налоги vs.\ долг) не влияет на совокупное потребление — значение имеет только приведённая стоимость будущих госрасходов. Эквивалентность нарушается при конечном горизонте жизни без альтруизма, финансовых ограничениях и искажающем налогообложении.
\item \textbf{Фискальные мультипликаторы}: эмпирически мультипликатор государственных расходов, как правило, выше мультипликатора снижения налогов, что согласуется с частичным сбережением дополнительного располагаемого дохода.
\item \textbf{Кривая Лаффера}: рост ставки налога на труд увеличивает поступления лишь до определённого предела, после которого искажение стимулов к труду начинает снижать налоговую базу быстрее, чем растёт ставка.
\end{itemize}

\newpage
\section{Лекция 3. Модель Рамсея}

\subsection{Мотивация}

В предыдущей лекции (Лекция 2) рассматривался межвременной выбор репрезентативного потребителя, однако трудовой доход и процентная ставка считались \textbf{экзогенными}. В более широкой картине эти переменные эндогенны и взаимосвязаны с выбором потребителей, поэтому требуется анализ общего равновесия.

Индивидуальные сбережения выбираются для перераспределения ресурсов во времени. Но что определяет \textbf{агрегированную} норму сбережений? Модель Рамсея ставит перед собой следующие вопросы:
\begin{itemize}[leftmargin=1.3em]
\item Что определяет оптимальную эндогенную агрегированную норму сбережений?
\item Можно ли обратиться к \textit{золотому правилу} накопления капитала?
\item Каковы эффекты фискальной политики в классической модели общего равновесия?
\item Ранее обсуждались различия взглядов кейнсианцев и классиков без микрооснований — модель Рамсея даёт эти микрооснования.
\end{itemize}

\subsection{Производство и накопление капитала}

Экономика состоит из бесконечно живущих домохозяйств. Население растёт с темпом $n$:
\begin{formulabox}
\begin{equation}
L_{t+1} = L_t(1+n) \tag{3.1$'$}
\end{equation}
\end{formulabox}

Каждое домохозяйство поставляет одну единицу труда. Выпуск $Y$ производится с помощью капитала $K$ и труда $L$. Производственная функция:
\begin{formulabox}
\begin{equation}
Y_t = F(K_t, L_t) \tag{3.1}
\end{equation}
\end{formulabox}
удовлетворяет неоклассическим предпосылкам:
$$F_K' > 0,\quad F_L' > 0,\quad F_{KK}'' < 0,\quad F_{LL}'' < 0,\quad F_{KL}'' > 0$$
(плюс линейная однородность и условия Инады).

Производственная функция в интенсивной форме (на одного работника):
\begin{formulabox}
\begin{equation}
y_t = f(k_t), \qquad y_t = Y_t/L_t,\quad k_t = K_t/L_t \tag{3.2}
\end{equation}
\end{formulabox}

Агрегатное тождество дохода:
\begin{formulabox}
\begin{equation}
Y_t = C_t + I_t \tag{3.3}
\end{equation}
\end{formulabox}

Инвестиции и накопление капитала:
\begin{formulabox}
\begin{align}
K_{t+1} &= I_t + K_t(1-\delta) \notag\\
&= Y_t - C_t + K_t(1-\delta) \tag{3.4}
\end{align}
\end{formulabox}

\subsubsection{Переход к подушевым переменным}

Разделим (3.4) на $L_t$:
$$\frac{K_{t+1}}{L_t} = \frac{Y_t}{L_t} - \frac{C_t}{L_t} + \frac{K_t}{L_t}(1-\delta), \qquad L_t = \frac{L_{t+1}}{1+n}.$$

Отсюда получаем ключевое \textbf{уравнение накопления капитала на одного работника}:
\begin{formulabox}
\begin{equation}
(1+n)k_{t+1} = f(k_t) - c_t + k_t(1-\delta) \tag{3.5}
\end{equation}
\end{formulabox}

\subsection{Накопление капитала и цены факторов}

При совершенной конкуренции цены факторов производства равны их предельным продуктам:
\begin{formulabox}
\begin{align}
w_t &= f(k_t) - k f'(k_t) \tag{3.6}\\
r_t &= f'(k_t) - \delta \tag{3.7}
\end{align}
\end{formulabox}

\begin{intuition}
Реальная заработная плата равна предельному продукту труда, а реальная ставка процента (доходность капитала) — предельному продукту капитала за вычетом нормы амортизации $\delta$, поскольку часть валового дохода от капитала идёт на возмещение его выбытия.
\end{intuition}

\subsection{Задача оптимизации}

Социальный планировщик максимизирует пожизненную полезность домохозяйств (общественное благосостояние) при ограничении на ресурсы:
\begin{formulabox}
\begin{equation}
\sum_{t=0}^{\infty} \frac{u(c_t)}{(1+\rho)^t} \to \max_{c_t} \tag{3.8}
\end{equation}
\end{formulabox}
при ограничении
\begin{formulabox}
\begin{equation*}
(1+n)k_{t+1} = f(k_t) - c_t + k_t(1-\delta) \tag{3.5}
\end{equation*}
\end{formulabox}

Позже решение планировщика будет сопоставлено с децентрализованным равновесием.

\textbf{Замечание.} Альтернативная (бентамианская) спецификация учитывает рост населения:
$$\sum_{t=0}^{\infty} \frac{L_t u(c_t)}{(1+\rho)^t} \to \max_{c_t} \quad \Leftrightarrow \quad \sum_{t=0}^{\infty} \frac{u(c_t)}{\big((1+\rho)/(1+n)\big)^t} \to \max_{c_t}, \qquad \rho > n$$

\subsubsection{Вывод условий первого порядка (Лагранжиан)}

Составим функцию Лагранжа:
$$\Lambda = \sum_{t=0}^{\infty} \frac{u(c_t)}{(1+\rho)^t} + \sum_{t=0}^{\infty} \lambda_t\Big(f(k_t) - c_t + k_t(1-\delta) - (1+n)k_{t+1}\Big).$$

Условие первого порядка по $c_t$:
$$\frac{\partial \Lambda}{\partial c_t} = 0 \quad \Rightarrow \quad \frac{u'(c_t)}{(1+\rho)^t} - \lambda_t = 0 \quad \forall t.$$

Условие первого порядка по $k_{t+1}$:
$$\frac{\partial \Lambda}{\partial k_{t+1}} = 0 \quad \Rightarrow \quad -\lambda_t(1+n) + \lambda_{t+1}\big(f'(k_{t+1}) + (1-\delta)\big) = 0 \quad \forall t.$$

Отсюда
$$\frac{\lambda_t}{\lambda_{t+1}} = \frac{1+f'(k_{t+1}) - \delta}{1+n} \quad \Rightarrow$$

\begin{formulabox}
\begin{equation}
\frac{u'(c_t)}{u'(c_{t+1})} = \frac{1+f'(k_{t+1})-\delta}{(1+n)(1+\rho)} \tag{3.9}
\end{equation}
\end{formulabox}

\subsection{Уравнение Эйлера}

Подставляя (3.7) в (3.9), получаем уравнение Эйлера для подушевого потребления:
\begin{formulabox}
\begin{equation}
\frac{u'(c_t)}{u'(c_{t+1})} = \frac{1+r_{t+1}}{(1+n)(1+\rho)} \tag{3.10}
\end{equation}
\end{formulabox}

Предположим изоэластичную (CRRA) функцию полезности:
\begin{formulabox}
\begin{equation}
u(c) = \frac{c^{1-1/\sigma}}{1-1/\sigma}, \tag{3.11}
\end{equation}
\end{formulabox}
где $\sigma$ — межвременная (мгновенная) эластичность замещения.

Тогда
$$\frac{c_{t+1}}{c_t} = \left(\frac{1+r_{t+1}}{(1+n)(1+\rho)}\right)^\sigma.$$

Линеаризуя (используя приближение $(1+x)^\sigma \approx 1+\sigma x$ при малых $x$):
$$\left(\frac{1+r_{t+1}}{(1+n)(1+\rho)}\right)^\sigma \approx 1 + \sigma\left(\frac{1+r_{t+1}}{(1+n)(1+\rho)} - 1\right) \approx 1 + \sigma(r_{t+1} - \rho - n)$$

Отсюда получаем приближённую динамику потребления в дискретном времени:
\begin{formulabox}
\begin{align}
\frac{c_{t+1}-c_t}{c_t} &\approx \sigma(r_{t+1} - \rho - n) \notag\\
&= \sigma\big(f'(k_{t+1}) - \rho - n - \delta\big) \tag{3.12}
\end{align}
\end{formulabox}

\begin{intuition}
Уравнение Эйлера связывает темп роста потребления с разницей между рыночной доходностью капитала $r_{t+1}$ и "эффективной" нормой дисконтирования $\rho+n$. Если доходность капитала превышает субъективную норму дисконта (с поправкой на рост населения), домохозяйству выгодно откладывать потребление на будущее — потребление растёт со временем ($\dot c>0$). Параметр $\sigma$ измеряет, насколько сильно потребитель готов менять профиль потребления в ответ на изменение ставки процента.
\end{intuition}

\subsection{Динамическая система}

Совместная динамика капитала и потребления:
\begin{formulabox}
\begin{align*}
(1+n)k_{t+1} &= f(k_t) - c_t + k_t(1-\delta), \tag{3.5}\\
\frac{c_{t+1}-c_t}{c_t} &= \sigma\big(f'(k_{t+1}) - \rho - n - \delta\big) \tag{3.12}
\end{align*}
\end{formulabox}

Для удобства перейдём к непрерывному времени:
\begin{formulabox}
\begin{align}
\dot{k}(t) &= f(k(t)) - c(t) - (n+\delta)k(t), \tag{3.13}\\
\frac{\dot{c}(t)}{c(t)} &= \sigma\big(f'(k(t)) - \rho - n - \delta\big) \tag{3.14}
\end{align}
\end{formulabox}

\subsection{Стационарное состояние (steady state)}

В стационарном состоянии переменные на одного работника не меняются во времени: $\dot k = 0$ и $\dot c = 0$.

Из $\dot k = 0$:
\begin{formulabox}
\begin{equation}
c^{*} = f(k^{*}) - (n+\delta)k^{*} \tag{3.15}
\end{equation}
\end{formulabox}

Из $\dot c = 0$:
\begin{formulabox}
\begin{equation}
f'(k^{*}) = \rho + n + \delta \tag{3.16}
\end{equation}
\end{formulabox}

\begin{intuition}
Уравнение (3.16) — это \textbf{модифицированное золотое правило}: в стационарном состоянии предельный продукт капитала равен сумме субъективной нормы дисконтирования, темпа роста населения и нормы амортизации. Это условие определяет уникальный уровень $k^{*}$, к которому сходится экономика, независимо от начального уровня капитала на разумных траекториях.
\end{intuition}

\subsection{Фазовая диаграмма}

Система (3.13)--(3.14) описывается фазовой диаграммой в координатах $(k,c)$:
\begin{itemize}[leftmargin=1.3em]
\item Локус $\dot k = 0$ — это кривая $c = f(k) - (n+\delta)k$ — "колоколообразная" кривая, возрастающая, затем убывающая (в форме кривой, как в модели Солоу); выше неё $\dot k < 0$, ниже — $\dot k > 0$.
\item Локус $\dot c = 0$ — это \textbf{вертикальная} прямая $k = k^{*}$, определяемая условием $f'(k^{*}) = \rho+n+\delta$; левее неё $f'(k) > \rho+n+\delta$, значит $\dot c > 0$ (потребление растёт), правее — $\dot c < 0$.
\end{itemize}

Пересечение этих двух локусов задаёт стационарную точку $(k^{*}, c^{*})$. Из анализа стрелок движения следует, что точка $(k^*,c^*)$ является \textbf{седловой} точкой: существует единственная сходящаяся (седловая) траектория, ведущая к стационарному состоянию из любого начального $k_0$; все остальные траектории либо уходят в бесконечность по $c$, либо приводят к $k=0$ с положительным $c$ (что невозможно поддерживать долго). Оптимальное поведение потребления в момент $t=0$ выбирается так, чтобы экономика сразу попала на седловую траекторию (например, при заданном $k_0$ выбирается $c_0$, соответствующее седловой траектории).

\subsection{Агрегированная норма сбережений в стационарном состоянии}

Введём норму сбережений: $s = 1 - C/Y = 1 - c/f(k)$. В этой модели норма сбережений \textbf{эндогенна и оптимальна}.

$$s^{*} = 1 - \frac{c^{*}}{f(k^{*})}$$

Используя $c^{*} = f(k^{*}) - (n+\delta)k^{*}$:
$$s^{*} = (n+\delta)\frac{k^{*}}{f(k^{*})} = (n+\delta)\frac{f'(k^{*})k^{*}}{f(k^{*})f'(k^{*})}.$$

Обозначим эластичность выпуска по капиталу $\alpha = f'(k^{*})k^{*}/f(k^{*})$ и предположим, что она постоянна (то есть производственная функция Кобба--Дугласа). Тогда, используя $f'(k^{*}) = \rho+n+\delta$:

\begin{formulabox}
\begin{equation}
s^{*} = \frac{(n+\delta)\alpha}{\rho+n+\delta} \tag{3.17}
\end{equation}
\end{formulabox}

то есть $s^{*} = s^{*}(\alpha, n, \delta, \rho)$, причём $\partial s^*/\partial\alpha>0$, $\partial s^*/\partial n>0$, $\partial s^*/\partial\delta>0$, $\partial s^*/\partial\rho<0$.

\subsubsection{Свойства}

\begin{itemize}[leftmargin=1.3em]
\item $s^{*} > 0$, если $n+\delta > 0$: сбережения нужны, если население растёт и/или капитал изнашивается. Чем быстрее экономический рост, тем выше норма сбережений.
\item Если $\rho = 0$, то $s^{*} = \alpha$ — это "классическое" золотое правило (как в модели Солоу).
\item Если $\rho > 0$, то $s^{*} < \alpha$ — нетерпеливость снижает норму сбережений. \textbf{Важно}: это не связано напрямую с уравнением Эйлера (в стационарном состоянии потребление постоянно). Более высокое $\rho$ означает более низкий $k^*$, а значит более низкие $f(k^*)$ и $c^*$:
$$c^{*} = f(k^{*}) - (n+\delta)k^{*}, \qquad f'(k^{*}) = \rho+n+\delta$$
$$\frac{\partial c^{*}}{\partial k^{*}} = f'(k^{*}) - (n+\delta) = \rho < f'(k^{*}) \quad \Rightarrow \quad \frac{\partial [s^{*}f(k^{*})]}{\partial k^{*}} < 0.$$
\end{itemize}

\subsection{Модифицированное золотое правило}

На фазовой диаграмме уровень $k^{gr}$ (обычное золотое правило) максимизирует потребление $c^{\max}$ и характеризуется условием $f'(k^{gr}) = n+\delta$ (без слагаемого $\rho$). Стационарное равновесие Рамсея $k^{*}$ определяется условием $f'(k^{*}) = \rho+n+\delta$.

Поскольку $\rho > 0$, а $f''(\cdot) < 0$ (убывающая отдача капитала):
$$k^{*} < k^{gr}, \qquad c^{*} < c^{\max}.$$

\begin{intuition}
Нетерпеливые домохозяйства ($\rho>0$) не желают накапливать капитал до уровня, максимизирующего стационарное потребление ($k^{gr}$), поскольку это требует чрезмерного откладывания текущего потребления. Оптимальный (с точки зрения благосостояния) уровень капитала $k^*$ меньше "золотого" уровня $k^{gr}$: планировщик учитывает нетерпеливость потребителей.
\end{intuition}

\subsection{Общественное благосостояние против децентрализованного равновесия}

Решение задачи социального планировщика подразумевает $k^{*}$ и $c^{*}$ меньше уровней золотого правила — они учитывают нетерпеливость агентов, но при этом являются \textbf{оптимизирующими благосостояние} уровнями.

Возникает вопрос: выберут ли децентрализованные домохозяйства то же самое? В чём разница между задачей планировщика и задачей децентрализованных домохозяйств?
\begin{itemize}[leftmargin=1.3em]
\item Функция полезности одинакова.
\item Но социальный планировщик учитывает ресурсное ограничение \textbf{всей экономики}, тогда как домохозяйства учитывают \textbf{свои собственные} бюджетные ограничения.
\end{itemize}

\subsubsection{Бюджетное ограничение домохозяйства}

Бюджетное ограничение домохозяйства:
\begin{formulabox}
\begin{equation}
C_t + (K_{t+1} - K_t) = w_t L_t + r_t K_t \tag{3.18}
\end{equation}
\end{formulabox}
(при совершенной конкуренции/условии нулевой прибыли неважно, кому принадлежат фирмы).

Разделим на $L_t$ и запишем в подушевых терминах:
\begin{formulabox}
\begin{equation}
c_t + (1+n)k_{t+1} - k_t = w_t + r_t k_t \tag{3.19}
\end{equation}
\end{formulabox}

Подставляя определения цен факторов (3.6)--(3.7), приходим к тому же самому ресурсному ограничению:
\begin{formulabox}
\begin{equation}
(1+n)k_{t+1} = f(k_t) - c_t + k_t(1-\delta) \tag{3.5}
\end{equation}
\end{formulabox}
\begin{intuition}
Динамическое бюджетное ограничение домохозяйств совпадает с ресурсным ограничением социального планировщика. Следовательно, децентрализованное равновесие (выбор потребления и накопления капитала) \textbf{Парето-эффективно} — потому что решение планировщика (так называемый \textit{командный оптимум}) Парето-эффективно по построению. Это означает, что \textbf{любое} возмущение в виде фискальной политики/перераспределения \textbf{не может} быть Парето-улучшающим.
\end{intuition}

\subsection{Фискальная политика при сбалансированном бюджете: паушальный (lump-sum) налог}

Пусть государство поддерживает сбалансированный бюджет и финансирует (бесполезные) расходы паушальным налогом:
\begin{formulabox}
\begin{equation}
g_t = \tau_t \tag{3.20}
\end{equation}
\end{formulabox}

В этом случае бюджетное ограничение домохозяйства снова идентично ресурсному ограничению планировщика:
\begin{formulabox}
\begin{align}
c_t + (1+n)k_{t+1} - k_t &= w_t + r_t k_t - \tau_t \tag{3.21}\\
(1+n)k_{t+1} &= f(k_t) - c_t - g_t + k_t(1-\delta) \tag{3.22}
\end{align}
\end{formulabox}

Уравнение Эйлера (3.10) остаётся неизменным: паушальный налог \textbf{не влияет} на цену межвременного выбора. Таким образом, локус $\dot c=0$ остаётся прежним, а локус $\dot k=0$ сдвигается вниз (на величину $g$).

\begin{intuition}
Возникает поразительный результат: государственные расходы, финансируемые паушальным налогом, \textbf{полностью вытесняют} частное потребление, не оказывая влияния ни на процентную ставку, ни на накопление капитала. В стационарном состоянии чистые инвестиции равны нулю как с вмешательством государства, так и без него. Стационарный капитал на одного работника не меняется, поэтому "безубыточные" инвестиции $(n+\delta)k$ остаются прежними, а всё сокращение ресурсов ложится на потребление домохозяйств один к одному.
\end{intuition}

\subsection{Фискальная политика при сбалансированном бюджете: искажающий налог на доход от капитала}

Пусть теперь государство облагает доход от капитала по ставке $\tau_K$ и перераспределяет весь доход в виде паушального трансферта $x$:
\begin{formulabox}
\begin{equation}
x_t = \tau_K f'(k_t) k_t \tag{3.23}
\end{equation}
\end{formulabox}

На первый взгляд кажется, что доход и, следовательно, бюджетное ограничение потребителя не меняются. Однако искажающая политика меняет посленалоговую процентную ставку, а значит и уравнение Эйлера:
\begin{formulabox}
\begin{align}
r_t &= (1-\tau_K) f'(k_t) - \delta \tag{3.24}\\
\frac{c_{t+1}-c_t}{c_t} &= \sigma\big((1-\tau_K) f'(k_{t+1}) - \rho - n - \delta\big) \tag{3.25}
\end{align}
\end{formulabox}

Локус $\dot k = 0$ остаётся прежним, а локус $\dot c = 0$ сдвигается \textbf{влево} (стационарный $k^{*}$ становится меньше).

\begin{intuition}
Перераспределительная политика вносит Парето-ухудшающие искажения. Хотя бюджетное ограничение домохозяйства как уравнение выглядит неизменным, политика заставляет домохозяйства изменить выбор между потреблением и сбережением. Более низкая посленалоговая ставка процента снижает стимулы к сбережению, что ведёт к меньшему накоплению капитала ($k^{*}\downarrow$), а значит к меньшему доходу ($f(k^{*})\downarrow$) и потреблению ($c^{*}\downarrow$).

Формально это действует так же, как если бы потребители стали более нетерпеливыми ($\rho^{*}\uparrow \to k^{*}\downarrow$). Но если нетерпеливые потребители потребляют меньше \textbf{оптимальным} образом, то перераспределение через искажающий налог \textbf{не является} Парето-эффективным.

Обратное перераспределение (субсидия на доход от капитала $\varsigma_K$, финансируемая паушальным налогом $\tau_t=\varsigma_K f'(k_t)k_t$) сдвигает локус $\dot c=0$ вправо и может достичь максимального $c^{*}$ (уровня золотого правила). Но такое принудительное перепотребление также \textbf{не является} оптимальным!
\end{intuition}

\newpage
\section{Лекции 4--5. Модель перекрывающихся поколений (OLG)}

\subsection{Мотивация}

Из лекции 3 известно, что децентрализованное равновесие в экономике с бесконечно живущими домохозяйствами Парето-эффективно, и были изучены свойства оптимальной эндогенной агрегированной нормы сбережений. Но универсален ли этот результат для всех динамических конкурентных экономик?

Теперь рассматривается другая конструкция — динамическая экономика \textbf{перекрывающихся поколений} конечно живущих домохозяйств. Это второй тип динамических моделей общего равновесия в макроэкономике.

Бесконечный горизонт жизни не позволяет изучать ряд важных вопросов:
\begin{itemize}[leftmargin=1.3em]
\item Важным компонентом богатства в любой экономике является \textbf{пенсионное богатство}. Для изучения этого вопроса необходимо разделить конечную жизнь на трудовой период и период выхода на пенсию.
\item Важны \textbf{межпоколенческие отношения}.
\item Будет показано, что модифицированное золотое правило \textbf{не обязательно} выполняется при конечном горизонте жизни.
\end{itemize}

\subsection{Структура модели: перекрывающиеся поколения}

Экономика состоит из перекрывающихся поколений. Каждое домохозяйство в поколении живёт \textbf{два периода}. В каждом периоде $t$ сосуществуют два поколения:
\begin{itemize}[leftmargin=1.3em]
\item \textbf{Молодое} поколение живёт в периодах $t$ и $t+1$;
\item \textbf{Старое} поколение живёт в периодах $t-1$ и $t$.
\end{itemize}

Размер поколений растёт с темпом $n$:
$$L_{t+1} = L_t(1+n).$$

\textbf{Молодые домохозяйства:}
\begin{itemize}[leftmargin=1.3em]
\item поставляют одну единицу труда и получают заработную плату $w$;
\item потребляют $c_1$ и сберегают $s$ (из заработной платы).
\end{itemize}

\textbf{Старые домохозяйства:}
\begin{itemize}[leftmargin=1.3em]
\item не работают и не получают заработную плату;
\item используют доход от капитала $(1+r)s$ для потребления $c_2$.
\end{itemize}

\subsection{Задача домохозяйства}

Домохозяйство, живущее в периодах $t$ и $t+1$, максимизирует:
\begin{formulabox}
\begin{equation}
U(c_{1t}, c_{2t+1}) = u(c_{1t}) + \frac{u(c_{2t+1})}{1+\rho} \to \max_{c_{1t}, c_{2t+1}} \tag{4.1}
\end{equation}
\end{formulabox}
при бюджетных ограничениях:
\begin{formulabox}
\begin{align}
c_{1t} + s_t &= w_t \tag{4.2}\\
c_{2t+1} &= (1+r_{t+1})s_t \tag{4.3}
\end{align}
\end{formulabox}
или, объединив в единое (межвременное) бюджетное ограничение:
\begin{formulabox}
\begin{equation}
\frac{c_{2t+1}}{1+r_{t+1}} + c_{1t} = w_t \tag{4.4}
\end{equation}
\end{formulabox}

Уравнение Эйлера:
\begin{formulabox}
\begin{equation}
\frac{u'(c_{1t})}{u'(c_{2t+1})} = \frac{1+r_{t+1}}{1+\rho} \tag{4.5}
\end{equation}
\end{formulabox}
или, подставив бюджетные ограничения явно:
\begin{formulabox}
\begin{equation}
\frac{u'(w_t - s_t)}{u'\big((1+r_{t+1})s_t\big)} = \frac{1+r_{t+1}}{1+\rho} \tag{4.6}
\end{equation}
\end{formulabox}

\subsection{Сбережения}

Из (4.6) следует, что $s$ является неявной функцией $w$ и $r$:
$$s_t = s_t(w_t, r_{t+1}).$$

Можно показать, что
$$0 < \frac{\partial s_t}{\partial w_t} < 1$$
(рост зарплаты увеличивает сбережения, но не на всю величину прироста — часть идёт на текущее потребление).

Предполагается, что эффект межвременного замещения доминирует над эффектом дохода:
$$\frac{\partial s_t}{\partial r_{t+1}} > 0$$

\begin{intuition}
Рост будущей ставки процента имеет два противоположных эффекта на сбережения молодого поколения: эффект замещения (более высокая доходность делает будущее потребление "дешевле" в терминах текущего сбережения, что стимулирует больше сберегать) и эффект дохода (более высокая доходность позволяет достичь того же будущего потребления с меньшими сбережениями). Предположение $\partial s_t/\partial r_{t+1}>0$ означает, что эффект замещения преобладает.
\end{intuition}

\subsection{Производство и цены факторов}

Производственный сектор такой же, как в модели Рамсея (и модели Солоу):
\begin{formulabox}
\begin{align}
Y_t &= F(K_t, L_t), \tag{4.7}\\
y_t &= f(k_t), \quad y_t = Y_t/L_t,\ k_t = K_t/L_t \tag{4.8}
\end{align}
\end{formulabox}

Цены факторов при совершенной конкуренции:
\begin{formulabox}
\begin{align}
w_t &= f(k_t) - k f'(k_t) \tag{4.9}\\
r_t &= f'(k_t) - \delta \tag{4.10}
\end{align}
\end{formulabox}

\subsection{Накопление капитала}

Предполагается, что инвестиции в периоде $t$ создают капитал для периода $t+1$, а старые домохозяйства продают свой капитал молодым, не оставляя наследства (bequest):
\begin{formulabox}
\begin{equation}
K_{t+1} = s_t L_t \tag{4.11}
\end{equation}
\end{formulabox}

Разделим на $L_{t+1}$:
\begin{formulabox}
\begin{align}
k_{t+1} &= \frac{s_t}{1+n} \notag\\
&= \frac{s_t(w_t, r_{t+1})}{1+n} \notag\\
&= \frac{s\big(f(k_t) - k_t f'(k_t),\ f'(k_{t+1}) - \delta\big)}{1+n} \tag{4.12}
\end{align}
\end{formulabox}

Уравнение (4.12) — это разностное уравнение первого порядка (с $k_{t+1}$ по обе стороны, поскольку $r_{t+1}$ зависит от $k_{t+1}$), задающее траекторию капитала во времени.

\subsection{Стационарное состояние}

Можно показать, что
$$0 < \dfrac{-\dfrac{\partial s_t}{\partial w_t}k^{*}f''(k^{*})}{1+n-\dfrac{\partial s_t}{\partial r_{t+1}}f''(k^{*})} < 1.$$

При определённых предположениях стационарное состояние \textbf{единственно и устойчиво} (это верно, например, для случая логарифмической функции полезности и производственной функции Кобба--Дугласа). Графически это отражается на диаграмме $k_{t+1}(k_t)$: вогнутая кривая пересекает биссектрису $k_{t+1}=k_t$ в единственной внутренней точке $k^{*}$, и итерационный процесс $k_0 \to k_1 \to \dots \to k^{*}$ сходится к ней.

\subsubsection{Уравнение Эйлера в стационарном состоянии}

Рассмотрим уравнение Эйлера (4.5) в стационарном состоянии:
$$\frac{u'(c_1^{*})}{u'(c_2^{*})} = \frac{1+r^{*}}{1+\rho}.$$

Если бы выполнялось $u'(c_1^{*})/u'(c_2^{*}) = 1+n$, то мы получили бы модифицированное золотое правило, как в модели Рамсея:
$$(1+r^{*}) = (1+\rho)(1+n) \ \Leftrightarrow\ r^{*} \approx \rho+n \ \Leftrightarrow\ f'(k^{*}) \approx \rho+n+\delta.$$

Однако \textbf{в общем случае это равенство не выполняется} — иначе капитал не накапливается оптимальным образом (см. далее раздел о динамической эффективности).

\subsection{Парето-эффективность: Рамсей против OLG}

\begin{itemize}[leftmargin=1.3em]
\item В модели Рамсея решение социального планировщика оптимально по построению, и динамическое бюджетное ограничение бесконечно живущих домохозяйств совпадает с ресурсным ограничением всей экономики. Поэтому децентрализованное равновесие Парето-эффективно.
\item В модели OLG бюджетные ограничения атомистических домохозяйств из \textbf{разных поколений не обязательно} формируют ресурсное ограничение экономики. Домохозяйства с конечным сроком жизни \textbf{не учитывают}, что происходит с накоплением капитала после их смерти.
\item Оптимальные индивидуальные решения \textbf{не обязательно} составляют Парето-эффективное равновесие.
\end{itemize}

Чтобы изучить этот вопрос, рассматривается \textbf{командный оптимум}: вновь предполагается, что социальный планировщик максимизирует общественное благосостояние при ресурсном ограничении экономики.

\subsection{Общественное благосостояние}

Определим общественное благосостояние как сумму полезностей всех текущих и будущих поколений:
\begin{formulabox}
\begin{align}
W &= u(c_{20}) + \sum_{t=0}^{\infty} \frac{U(c_{1t}, c_{2t+1})}{(1+\rho)^t} \notag\\
&= u(c_{20}) + \sum_{t=0}^{\infty} \frac{u(c_{1t}) + \dfrac{u(c_{2t+1})}{1+\rho}}{(1+\rho)^t} \to \max_{c_{1t}, c_{2t}} \tag{4.13}
\end{align}
\end{formulabox}

Выбор функции общественного благосостояния подчинён разным взглядам:
\begin{itemize}[leftmargin=1.3em]
\item предполагается, что социальный планировщик \textbf{не учитывает} размер поколений (как и в модели Рамсея);
\item также предполагается, что планировщик использует ту же субъективную норму дисконтирования, что и домохозяйства (это на самом деле не принципиально).
\end{itemize}

\subsection{Ресурсное ограничение}

Социальный планировщик сталкивается с ресурсным ограничением всей экономики:
\begin{formulabox}
\begin{equation}
(1-\delta)K_t + F(K_t, L_t) = K_{t+1} + L_{t-1}c_{2t} + L_t c_{1t} \tag{4.14}
\end{equation}
\end{formulabox}

Разделим на $L_t$:
\begin{formulabox}
\begin{equation}
(1-\delta)k_t + f(k_t) = (1+n)k_{t+1} + c_{1t} + \frac{c_{2t}}{1+n} \tag{4.15}
\end{equation}
\end{formulabox}

Если определить агрегированное подушевое потребление как
$$c_t = c_{1t} + \frac{c_{2t}}{1+n},$$
то получаем \textbf{то же самое} ресурсное ограничение, что и в модели Рамсея:
\begin{formulabox}
\begin{equation}
(1-\delta)k_t + f(k_t) = (1+n)k_{t+1} + c_t \tag{4.15$'$}
\end{equation}
\end{formulabox}

\subsection{Командный оптимум}

Задача социального планировщика:
$$W = u(c_{20}) + \sum_{t=0}^{\infty} \frac{u(c_{1t}) + \dfrac{u(c_{2t+1})}{1+\rho}}{(1+\rho)^t} \to \max_{c_{1t}, c_{2t}}$$
$$\text{при} \quad (1-\delta)k_t + f(k_t) = (1+n)k_{t+1} + c_{1t} + \frac{c_{2t}}{1+n}$$

Лагранжиан:
$$\Lambda(c_1,c_2,k,\lambda) = u(c_{20}) + \sum_{t=0}^{\infty} \frac{u(c_{1t}) + \dfrac{u(c_{2t+1})}{1+\rho}}{(1+\rho)^t} + \sum_{t=0}^{\infty}\lambda_t\left[(1-\delta)k_t + f(k_t) - (1+n)k_{t+1} - c_{1t} - \frac{c_{2t}}{1+n}\right]$$

Условия первого порядка по $c_{1t}$, $c_{2t}$ и $k_{t+1}$ (проверяется прямым дифференцированием):
\begin{formulabox}
\begin{align}
\frac{u'(c_{1t})}{(1+\rho)^t} - \lambda_t &= 0 \tag{4.16}\\
\frac{u'(c_{2t})}{(1+\rho)^t} - \frac{\lambda_t}{1+n} &= 0 \tag{4.17}\\
\lambda_{t+1}\big(1+f'(k_{t+1}) - \delta\big) - \lambda_t(1+n) &= 0 \tag{4.18}
\end{align}
\end{formulabox}

Комбинируя (4.16)--(4.18) и используя (4.10) для процентной ставки, получаем:
\begin{formulabox}
\begin{align}
\frac{u'(c_{1t})}{u'(c_{2t})} &= 1+n \tag{4.19}\\
\frac{u'(c_{it})}{u'(c_{it+1})} &= \frac{1+r_{t+1}}{(1+\rho)(1+n)}, \quad i=1,2 \tag{4.20}
\end{align}
\end{formulabox}

\begin{intuition}
\textbf{Условие (4.19)} определяет оптимальное \textit{внутрипериодное} (intratemporal) распределение ресурсов между сосуществующими поколениями: предельная норма замещения между потреблением молодых и старых равна валовому темпу роста поколений $1+n$. Оптимальное распределение учитывает относительный размер поколений — если поколения растут быстрее, "вес" молодых в потреблении должен быть выше.

\textbf{Условие (4.20)} (аналогичное для молодых и старых агентов) определяет оптимальное \textit{межвременное} распределение ресурсов: предельная норма замещения между потреблением в последовательных периодах равна относительной цене $(1+r_{t+1})^{-1}$ с учётом субъективного дисконтирования будущей полезности и роста населения.
\end{intuition}

Комбинируя (4.19) и (4.20) при $i=2$, получаем \textbf{индивидуальное} уравнение Эйлера:
\begin{formulabox}
\begin{equation}
\frac{u'(c_{1t})}{u'(c_{2t+1})} = \frac{1+r_{t+1}}{1+\rho} \tag{4.5}
\end{equation}
\end{formulabox}

Итак, социальный планировщик выбирает распределение по тому же \textbf{принципу}, что и децентрализованные домохозяйства. Но означает ли это, что выбор совпадает? \textbf{Нет!}

\subsubsection{Почему выбор планировщика и домохозяйств различается}

Уравнения (4.5) и (4.20) — разные:
\begin{itemize}[leftmargin=1.3em]
\item Уравнение (4.5) определяет оптимальное межвременное распределение для \textbf{конкретного} домохозяйства, принимающего доход как \textbf{экзогенный} ($w_t$ зависит от накопления капитала в прошлом) и \textbf{игнорирующего} влияние своего выбора на доходы в будущем.
\item Уравнение (4.20) определяет предельную норму замещения между потреблением молодого (или старого) домохозяйства в периоде $t$ и потреблением молодого (или старого) из \textbf{следующего} поколения в периоде $t+1$. Перераспределение потребления между текущими и будущими молодыми (и старыми) домохозяйствами контролирует накопление капитала и означает \textbf{эндогенное} распределение дохода во времени — то, что отдельное домохозяйство не может и не пытается контролировать.
\end{itemize}

\subsection{Стационарное состояние командного оптимума}

В стационарном состоянии:
\begin{formulabox}
\begin{align}
\frac{u'(c_1^{*})}{u'(c_2^{*})} &= 1+n \tag{4.21}\\
1 &= \frac{1+r^{*}}{(1+\rho)(1+n)} \tag{4.22}
\end{align}
\end{formulabox}

Из (4.22):
$$(1+r^{*}) = (1+\rho)(1+n) \ \Leftrightarrow\ r^{*} \approx \rho+n \ \Leftrightarrow\ f'(k^{*}) \approx \rho+n+\delta.$$

\begin{intuition}
Таким образом, \textbf{командный оптимум} в модели OLG удовлетворяет \textbf{тому же самому} модифицированному золотому правилу накопления капитала, что и модель Рамсея, — в то время как децентрализованное равновесие с перекрывающимися поколениями \textbf{может ему не удовлетворять}. Это и есть ключевое расхождение между двумя типами моделей.
\end{intuition}

\subsection{Динамическая (не)эффективность}

Рассмотрим ресурсное ограничение в стационарном состоянии:
$$(1-\delta)k^{*} + f(k^{*}) = (1+n)k^{*} + c_1^{*} + \frac{c_2^{*}}{1+n}.$$

Введём стационарное агрегированное потребление $c^{*} = c_1^{*} + c_2^{*}/(1+n)$:
\begin{formulabox}
\begin{equation}
f(k^{*}) = c^{*} + (n+\delta)k^{*}, \qquad c^{*} = f(k^{*}) - (n+\delta)k^{*} \tag{4.23}
\end{equation}
\end{formulabox}

Рассмотрим изменение стационарного потребления, если социальный планировщик уменьшает стационарный капитал на $dk^{*} < 0$:
$$dc^{*} = \big(f'(k^{*}) - \delta - n\big)dk^{*} > 0 \quad \text{если} \quad f'(k^{*}) - \delta - n < 0,\ \text{т.е.}\ r^{*} < n.$$

\begin{intuition}
Если $r^{*} < n$ (доходность капитала ниже темпа роста населения/экономики), то стационарное состояние \textbf{Парето-неэффективно}: "проедая" часть капитала, можно \textbf{увеличить} агрегированное потребление, а значит и стационарное потребление \textbf{каждого} поколения. Это удивительный результат: при избыточном накоплении капитала (сверхнакоплении) сокращение капитала строго улучшает благосостояние всех — без каких-либо трансфертов между поколениями.
\end{intuition}

Тот же результат применим и к переходной динамике. Пусть планировщик уменьшает капитал во все будущие периоды: $dk_{t+\tau} = dk < 0,\ \tau>0$. Из ресурсного ограничения $(1-\delta)k_t+f(k_t)=(1+n)k_{t+1}+c_t$ получаем:
$$dc_t = -(1+n)dk > 0, \qquad dc_{t+\tau} = \big(f'(k_{t+\tau}) - n - \delta\big)dk > 0 \quad \text{если} \quad r_{t+\tau} < n.$$

Такая экономика называется \textbf{динамически неэффективной}. Это явление \textbf{динамическое} и потому принципиально отличается от микроэкономических провалов рынка (market failures).

\subsection{Альтруистические межпоколенческие трансферты}

Итак, с конечным сроком жизни возникают как минимум две "проблемы":
\begin{enumerate}[leftmargin=1.5em]
\item экономика может быть \textbf{динамически неэффективной};
\item \textbf{рикардианская эквивалентность} не выполняется, если перераспределение налогового бремени затрагивает разные поколения.
\end{enumerate}

Межпоколенческие трансферты / альтруистическое наследование могут решить эти проблемы. Предположим, что родители заботятся о полезности своих детей — тогда полезность поколения, живущего в периодах $t$ и $t+1$, принимает форму:
\begin{formulabox}
\begin{equation}
V_t = u(c_{1t}) + \frac{u(c_{2t+1})}{1+\rho} + \frac{V_{t+1}}{1+\theta} \to \max_{c_{1t},c_{2t+1}} \tag{4.24}
\end{equation}
\end{formulabox}

Для упрощения предположим, что родители дисконтируют полезность детей с той же ставкой, что и собственную будущую полезность: $\theta = \rho$.

\subsubsection{Бюджетные ограничения с трансфертами}

Поскольку население растёт с темпом $n$, у каждого репрезентативного домохозяйства $(1+n)$ "детей". Обозначим $x_t$ — трансферт (наследство) на одного "ребёнка". Тогда бюджетные ограничения:
\begin{formulabox}
\begin{align}
c_{1t} + s_t &= w_t + x_t \tag{4.25}\\
c_{2t+1} + (1+n)x_{t+1} &= (1+r_{t+1})s_t \tag{4.26}
\end{align}
\end{formulabox}
или
\begin{formulabox}
\begin{equation}
c_{1t} + \frac{c_{2t+1} + (1+n)x_{t+1}}{1+r_{t+1}} = w_t + x_t \tag{4.27}
\end{equation}
\end{formulabox}

Лагранжиан задачи:
$$\Lambda(c_{1t}, c_{2t+1}, x_{t+1}, \lambda_t) = \sum_{\tau=0}^{\infty} \frac{u(c_{1t+\tau}) + \dfrac{u(c_{2t+\tau+1})}{1+\rho}}{(1+\rho)^\tau} + \lambda_t\left[w_t + x_t - c_{1t} - \frac{c_{2t+1}+(1+n)x_{t+1}}{1+r_{t+1}}\right]$$

Условия первого порядка по $c_{1t}, c_{2t+1}, x_{t+1}$ (предполагая внутреннее решение):
\begin{formulabox}
\begin{align}
u'(c_{1t}) &= \lambda_t \tag{4.28}\\
\frac{u'(c_{2t+1})}{1+\rho} &= \frac{\lambda_t}{1+r_{t+1}} \tag{4.29}\\
\frac{u'(c_{1t+1})}{1+\rho} &= \frac{\lambda_t(1+n)}{1+r_{t+1}} \tag{4.30}
\end{align}
\end{formulabox}

\subsubsection{Стандартное уравнение Эйлера сохраняется}

Комбинируя (4.28) и (4.29), получаем стандартное уравнение Эйлера:
\begin{formulabox}
\begin{equation}
\frac{u'(c_{1t})}{u'(c_{2t+1})} = \frac{1+r_{t+1}}{1+\rho} \tag{4.5}
\end{equation}
\end{formulabox}

\begin{intuition}
Мотив наследования \textbf{не меняет} способ распределения личного потребления. Действительно, здесь есть две "раздельные" задачи: (1) выбор совокупного личного потребления и размера наследства; и (2) выбор распределения совокупного личного потребления между двумя периодами жизни.
\end{intuition}

\subsubsection{Межпоколенческие условия}

Из (4.28) и (4.30):
\begin{formulabox}
\begin{equation}
\frac{u'(c_{1t+1})}{u'(c_{2t+1})} = 1+n \tag{4.31}
\end{equation}
\end{formulabox}

Из (4.29) и (4.30):
\begin{formulabox}
\begin{equation}
\frac{u'(c_{1t+1})}{u'(c_{1t})} = \frac{(1+n)(1+\rho)}{1+r_{t+1}} \tag{4.32}
\end{equation}
\end{formulabox}

Аналогично:
\begin{formulabox}
\begin{equation}
\frac{u'(c_{2t+1})}{u'(c_{2t})} = \frac{(1+n)(1+\rho)}{1+r_{t+1}} \tag{4.33}
\end{equation}
\end{formulabox}

\subsubsection{Совпадение с командным оптимумом}

Условия (4.31)--(4.33) \textbf{идентичны} условиям (4.19)--(4.20, $i=1,2$) из задачи социального планировщика. Следовательно, \textbf{модифицированное золотое правило} $(1+r^{*}) = (1+\rho)(1+n)$ выполняется и здесь.

\begin{intuition}
Альтруистические домохозяйства выполняют ту же работу, что и социальный планировщик: они заботятся о будущих поколениях и, следовательно, \textbf{не} принимают свои доходы как экзогенные и \textbf{не} выбирают личное потребление эгоистически.

Ещё одно важное следствие: конечный срок жизни \textbf{не должен} разрушать рикардианскую эквивалентность, если домохозяйства оставляют альтруистическое наследство.

Можно легко ввести \textbf{двустороннний} альтруизм (когда и дети заботятся о родителях), однако его анализ значительно сложнее из-за "зеркального эффекта". Удивительно, но двусторонний альтруизм \textbf{не гарантирует} динамическую эффективность.
\end{intuition}

\subsection{Эффективный мотив наследования (Barro, 1974)}

Приведённые результаты \textbf{не являются общими}:
\begin{itemize}[leftmargin=1.3em]
\item Речь шла об \textbf{альтруистическом} наследовании. Наследство, основанное на иных мотивах, не приводит к тем же результатам!
\item Предполагалось внутреннее решение $x_{t+1}>0$ (и $x^{*}>0$). Барро (1974) называет это \textbf{эффективным мотивом наследования} (\textit{effective bequest motive}).
\item Но возможно и \textbf{угловое} решение $x_{t+1}=0$ (и $x^{*}=0$). Это может произойти, если:
  \begin{itemize}[leftmargin=1.3em]
  \item $n$ достаточно велико: иметь слишком много детей означает, что едва ли можно повысить их благосостояние за счёт сокращения собственного потребления;
  \item или $\theta \gg \rho$, то есть альтруизм родителей умеренный;
  \item или высокий рост производительности делает заработную плату будущих поколений достаточно выше, так что "бедным" родителям не нужно помогать своим "богатым" детям.
  \end{itemize}
\end{itemize}

\subsubsection{Угловое решение и динамическая неэффективность}

Угловое решение $x_{t+1}=0$ (и $x^{*}=0$) означает, что экономика \textbf{динамически неэффективна}:
\begin{formulabox}
\begin{align}
\frac{u'(c_{1t+1})}{u'(c_{2t+1})} &\le 1+n \tag{4.31$'$}\\
\frac{u'(c_{1t+1})}{u'(c_{1t})} &\le \frac{(1+n)(1+\rho)}{1+r_{t+1}} \tag{4.32$'$}\\
\frac{u'(c_{2t+1})}{u'(c_{2t})} &\le \frac{(1+n)(1+\rho)}{1+r_{t+1}} \tag{4.33$'$}
\end{align}
\end{formulabox}
$$\Rightarrow \quad (1+r^{*}) \le (1+\rho)(1+n).$$

\subsection{Системы социального обеспечения (пенсионные системы)}

До сих пор предполагалось, что домохозяйства сами заботятся о потреблении после выхода на пенсию. В реальном мире почти все экономики имеют "сети безопасности" (safety nets).
\begin{itemize}[leftmargin=1.3em]
\item Первая пенсионная система была введена в \textbf{1895 году Отто фон Бисмарком}.
\item Сегодня существует множество разных пенсионных систем.
\item Социальное обеспечение — более широкое понятие, включающее медицинское страхование, пособия по безработице и т.д.
\end{itemize}

С теоретической точки зрения можно выделить два типа пенсионных систем:
\begin{itemize}[leftmargin=1.3em]
\item \textbf{распределительная система} ("Pay as You Go", PAYG) — текущие работники финансируют текущих пенсионеров;
\item \textbf{накопительная система} (Fully-Funded) — работники формируют собственные пенсионные сбережения.
\end{itemize}

\subsubsection{Накопительная (Fully-Funded) пенсионная система}

Предположим, что молодое домохозяйство в периоде $t$ обязано внести взнос $m_t$ на пенсионный сберегательный депозит. Эта сумма инвестируется в производственный капитал. Таким образом, в следующем периоде $t+1$ домохозяйство получает пенсионную выплату $z_{t+1} = (1+r_{t+1})m_t$.

Предполагается, что $m_t$ не превышает уровня добровольных сбережений, которые выбирались бы без пенсионной системы.

\textbf{Вопрос:} влияет ли эта система на накопление капитала и равновесие?

\textbf{Ответ: нет!} Пенсионная система действует как \textbf{обязательное сбережение} в активе с той же доходностью:
\begin{formulabox}
\begin{align}
\frac{u'\big(w_t - (s_t+m_t)\big)}{u'\big((1+r_{t+1})(s_t+m_t)\big)} &= \frac{1+r_{t+1}}{1+\rho} \tag{4.34}\\
k_{t+1} &= \frac{s_t+m_t}{1+n} \tag{4.35}
\end{align}
\end{formulabox}

\begin{intuition}
Поскольку норма доходности пенсионного счёта совпадает с рыночной доходностью капитала $r$, домохозяйство просто снижает добровольные сбережения $s_t$ ровно на величину обязательного взноса $m_t$ (при внутреннем решении), оставляя совокупные сбережения $s_t+m_t$, а значит и накопление капитала, неизменными. Накопительная пенсионная система нейтральна.
\end{intuition}

\subsubsection{Распределительная (PAYG) пенсионная система}

Предположим, что молодое домохозяйство в периоде $t$ обязано внести взнос $m_t$ в пенсионную систему. Этот взнос передаётся старым домохозяйствам \textbf{в том же} периоде $t$. С учётом относительного размера поколений старое домохозяйство в периоде $t$ получает выплату $z_t = (1+n)m_t$.

Теперь пенсионная система \textbf{влияет} на накопление капитала и равновесие:
\begin{formulabox}
\begin{align}
\frac{u'\big(w_t - (s_t+m_t)\big)}{u'\big((1+r_{t+1})s_t + (1+n)m_{t+1}\big)} &= \frac{1+r_{t+1}}{1+\rho} \tag{4.36}\\
k_{t+1} &= \frac{s_t}{1+n} = \frac{s\big(w(k_t), r(k_{t+1}), m_t\big)}{1+n} \tag{4.37}
\end{align}
\end{formulabox}

\begin{intuition}
Уравнение Эйлера (4.36) отличается от уравнения (4.5): пенсионный взнос действует как обязательное сбережение в активе с нормой доходности $n \ne r$. Теперь у домохозяйств фактически есть "портфель" из двух разных активов — рыночного капитала (доходность $r$) и распределительной пенсионной системы (доходность $n$).
\end{intuition}

Как это влияет на добровольные сбережения?
$$n > r \ \Rightarrow\ \frac{\partial s_t}{\partial m_t} < -1, \qquad n < r \ \Rightarrow\ -1 < \frac{\partial s_t}{\partial m_t} < 0.$$

\begin{intuition}
Интуиция прямая: если $n>r$, то "пенсионный портфель" имеет доходность выше, чем доходность капитала, поэтому домохозяйство сокращает добровольные сбережения даже сильнее, чем на единицу возросшего взноса $m_t$ (замещение более чем один к одному, компенсируя выгодность обязательного "актива"); если $n<r$ — наоборот, замещение неполное.
\end{intuition}

Влияние на капитал:
\begin{formulabox}
\begin{equation}
\frac{dk_{t+1}}{dm_t} = \frac{\partial s_t/\partial m_t}{1+n-f''(k_{t+1})\,\partial s_t/\partial r_{t+1}} < 0 \tag{4.37$'$}
\end{equation}
\end{formulabox}

\begin{intuition}
Распределительная (PAYG) пенсионная система \textbf{изымает ресурсы} из накопления капитала: независимо от сбережений, инвестиции в производственный капитал становятся ниже, и стационарное отношение капитала к труду снижается.

Но всегда ли это плохо? \textbf{Нет.} Это может быть \textbf{Парето-улучшением}, если экономика была \textbf{динамически неэффективной}. В экономике с \textbf{эффективными мотивами наследования} PAYG-система действует как система паушальных налогов и трансфертов, что соответствует рикардианской эквивалентности. Однако на практике паушальные налоги встречаются редко.
\end{intuition}

\subsection{Пенсионные системы: дополнительные соображения}

Влияет ли выбор пенсионной системы на результат? Почему большинство стран столетие назад начинали с систем PAYG и сейчас пытаются перейти к накопительным системам?

Простая модель OLG не учитывает ряд важных аспектов реальности:
\begin{itemize}[leftmargin=1.3em]
\item \textbf{Старение населения.}
  \begin{itemize}[leftmargin=1.3em]
  \item Более низкое $n$ означает более низкую доходность в системе PAYG.
  \item Но также включает более высокую ожидаемую продолжительность жизни, что удлиняет период пенсии.
  \end{itemize}
\item \textbf{Эндогенный выбор возраста выхода на пенсию} — также вопрос политического выбора.
\item \textbf{Неполнота финансовых рынков} — рискованные инвестиции оставляют будущих пенсионеров незащищёнными.
\end{itemize}

Выбор между системой с фиксированными взносами и системой с фиксированными выплатами должен быть частью устойчивой фискальной политики.

\subsubsection{Реформы, экспроприация или финансовые репрессии}

\begin{itemize}[leftmargin=1.3em]
\item Если у страны система PAYG и стареющее население — лучше перейти к накопительной системе. Иначе придётся либо сокращать пенсионные выплаты/повышать пенсионный возраст, либо финансировать дефицит пенсионной системы через налоги и государственный долг, что порождает свою проблему.
\item Но \textbf{кто заплатит за реформу}?
\item Если уже есть смешанная система, и в её распределительной части дефицит — возникает соблазн использовать пенсионные накопления \textbf{будущих} пенсионеров для финансирования выплат \textbf{текущим} пенсионерам (тем, кто начинал ещё при системе PAYG). Из новейшей истории экспроприации: Португалия, Венгрия (2011), Польша (2014), Россия (...).
\item Если в стране накопительная система, а бремя государственного долга велико и власти опасаются "bond vigilantes" (инвесторов, требующих премию за риск по гособлигациям) — существует соблазн убедить/обязать пенсионные фонды покупать только государственные облигации (финансовые репрессии).
\end{itemize}

\subsection{Иллюстративные данные: пенсионные фонды в мире (OECD, 2014)}

Слайды лекции иллюстрируют эти теоретические выводы эмпирическими данными OECD "Pension Markets in Focus, 2014":
\begin{itemize}[leftmargin=1.3em]
\item \textbf{Размер пенсионных активов как \% ВВП} сильно различается между странами: от свыше 100\% ВВП (Нидерланды — 166{,}3\%, Исландия — 148{,}7\%, Швейцария — 119{,}0\%, Австралия — 103{,}3\%, Великобритания — 100{,}7\%) до долей процента ВВП (Франция — 0{,}4\%, Греция — 0{,}1\%); среди развивающихся стран лидируют ЮАР (87{,}1\%), Намибия (76{,}6\%), Лихтенштейн (65{,}8\%).
\item \textbf{Средняя доходность пенсионных активов за 5 лет} (номинальная/реальная, \%): например, Нидерланды 9{,}6/7{,}4, Канада 9{,}1/7{,}4, США 7{,}9/5{,}7, Греция 1{,}5/$-0{,}3$; наибольшая номинальная доходность — Пакистан 14{,}0/3{,}2, Колумбия 13{,}3/10{,}4.
\item \textbf{Структура активов пенсионных фондов} (доли акций, облигаций, депозитов и прочего) также сильно различается: например, в США доля акций около 49\%, тогда как в Корее — менее 10\%, с преобладанием депозитов.
\end{itemize}

\begin{intuition}
Эти эмпирические различия иллюстрируют, что теоретическая модель — лишь упрощённый каркас: реальные пенсионные системы различаются по масштабу, структуре активов и доходности, что отражает разные институциональные traditions, уровень развития финансовых рынков и политику регулирования в разных странах.
\end{intuition}

\subsection{Краткое резюме ключевых результатов (Лекции 3--5)}

\begin{itemize}[leftmargin=1.3em]
\item \textbf{Модель Рамсея} (бесконечно живущие домохозяйства): децентрализованное равновесие \textbf{всегда} Парето-эффективно; стационарное состояние удовлетворяет \textbf{модифицированному золотому правилу} $f'(k^{*})=\rho+n+\delta$, причём $k^{*}<k^{gr}$ из-за нетерпеливости; норма сбережений эндогенна: $s^{*}=(n+\delta)\alpha/(\rho+n+\delta)$; паушальный налог полностью вытесняет частное потребление без искажений; искажающий налог на капитал сдвигает экономику от оптимума.
\item \textbf{Модель OLG} (конечно живущие домохозяйства, перекрывающиеся поколения): децентрализованное равновесие \textbf{не обязательно} Парето-эффективно — возможна \textbf{динамическая неэффективность} ($r^{*}<n$); командный оптимум восстанавливает то же модифицированное золотое правило, что и в модели Рамсея; \textbf{эффективный альтруистический мотив наследования} восстанавливает эффективность и рикардианскую эквивалентность; распределительная (PAYG) пенсионная система снижает накопление капитала, но может быть Парето-улучшающей именно в случае динамически неэффективной экономики; накопительная пенсионная система нейтральна по отношению к равновесию.
\end{itemize}

\newpage
\section{Лекция 6. Гипотеза жизненного цикла (Life-Cycle Hypothesis)}

\subsection{Мотивация}

Модели Рамсея и OLG дают широкую картину потребления, сбережений, инвестиций и накопления капитала. Но чтобы построить даже простую модель общего равновесия, приходится опускать некоторые детали и потенциально важные вопросы. В этих моделях сбережения напрямую направлялись в инвестиции.

Необходимо вернуться к рассмотрению \textbf{частичного равновесия}, чтобы:
\begin{itemize}[leftmargin=1.3em]
\item подробно изучить решения домохозяйств о потреблении и сбережениях;
\item исследовать, как фирмы (а не домохозяйства) выбирают инвестиции.
\end{itemize}
После этого будут рассмотрены финансовые рынки и финансовое посредничество.

\subsection{Структура модели жизненного цикла (LCH)}

Рассмотрим репрезентативное домохозяйство, которое входит в экономику в периоде $t$. Его горизонт жизни длится $T$ периодов (лет) и делится на два под-периода — трудовой и пенсионный:
$$T = L + R.$$

Трудовой доход (поток дохода от "надела") положителен и постоянен в течение $L$ трудовых периодов и равен нулю после выхода на пенсию:
\begin{formulabox}
\begin{align}
y_{t+\tau}^{t} &= y^{t}, \qquad 0 \le \tau \le L-1, \notag\\
y_{t+\tau}^{t} &= 0, \qquad L \le \tau \le T-1 \tag{6.$0$}
\end{align}
\end{formulabox}
где верхний индекс $t$ обозначает "год рождения" (когорту), $\tau$ — "возраст" (расстояние от рождения), а нижний индекс $t+\tau$ обозначает текущий календарный период.

\subsection{Задача домохозяйства}

Домохозяйство максимизирует стандартную пожизненную полезность:
\begin{formulabox}
\begin{equation}
\sum_{\tau=0}^{T-1} \frac{u(c_{t+\tau}^{t})}{(1+\rho)^\tau} \to \max_{c_{t+\tau}^{t}} \tag{6.1}
\end{equation}
\end{formulabox}

Для упрощения изложения предполагается, что процентная ставка и субъективная норма дисконтирования равны нулю: $r=\rho=0$. Тогда межвременное бюджетное ограничение принимает простой вид:
\begin{formulabox}
\begin{equation}
\sum_{\tau=0}^{T-1} c_{t+\tau}^{t} = \sum_{\tau=0}^{T-1} y_{t+\tau}^{t} \tag{6.2}
\end{equation}
\end{formulabox}

Тогда уравнение Эйлера предполагает \textbf{постоянный профиль потребления}:
\begin{formulabox}
\begin{equation}
\frac{u'(c_{t+\tau}^{t})}{u'(c_{t+\tau+1}^{t})} = \frac{1+r}{1+\rho}, \qquad r=\rho=0 \ \Rightarrow\ c_{t+\tau}^{t} = c_{t+\tau+1}^{t} = c^{t} \quad \forall \tau \tag{6.3}
\end{equation}
\end{formulabox}

\begin{intuition}
При $r=\rho$ уравнение Эйлера сводится к простому равенству предельных полезностей в соседних периодах, а значит — к постоянству самого потребления во времени. Это частный случай общего результата Лекции 2: динамика потребления определяется исключительно соотношением $r$ и $\rho$, а не профилем дохода.
\end{intuition}

\subsection{Индивидуальная предельная склонность к потреблению}

Учитывая предположение о трудовом доходе на протяжении жизненного цикла, из (6.2) и (6.3):
$$Tc^{t} = Ly^{t} + R\cdot 0,$$
откуда
\begin{formulabox}
\begin{equation}
c^{t} = \frac{L}{L+R}\,y^{t} \tag{6.4}
\end{equation}
\end{formulabox}

\begin{intuition}
Домохозяйство "размазывает" совокупный трудовой доход, заработанный за $L$ трудовых периодов, равномерно на все $T=L+R$ периода жизни. Коэффициент $L/(L+R)$ — это доля дохода, которую можно потратить в каждом отдельном периоде, если сглаживать потребление на весь жизненный цикл, включая пенсию.
\end{intuition}

\subsection{Сбережения на протяжении трудового периода}

Определим сбережения обычным образом: $s_{t+\tau}^{t} = y_{t+\tau}^{t} - c_{t+\tau}^{t}$. В течение трудового периода ($0\le\tau\le L-1$), поскольку $y_{t+\tau}^{t}=y^{t}$:
\begin{formulabox}
\begin{equation}
s_{t+\tau}^{t} = \frac{R}{L+R}\,y^{t}, \qquad 0 \le \tau \le L-1 \tag{6.5}
\end{equation}
\end{formulabox}

Для правдоподобной параметризации $L=40$ и $R=10$ получаем предельную склонность к потреблению $0{,}8$ и предельную склонность к сбережению $0{,}2$.

\subsection{Накопление и истощение активов}

Активы растут на протяжении трудового периода (начиная с нуля, поскольку начальное богатство и наследство не вводятся):
\begin{formulabox}
\begin{equation}
a_{t+\tau}^{t} = \tau s_{t+\tau}^{t} = \frac{\tau R}{L+R}\,y^{t}, \qquad 0 \le \tau \le L \tag{6.6}
\end{equation}
\end{formulabox}

и достигают максимума в последнем периоде перед выходом на пенсию:
\begin{formulabox}
\begin{equation}
a_{t+L}^{t} = \frac{LR}{L+R}\,y^{t}, \tag{6.7}
\end{equation}
\end{formulabox}
что при $L=40,\ R=10$ в \textbf{8 раз} превышает годовой доход.

После выхода на пенсию домохозяйство начинает жить за счёт накопленного богатства ($y_{t+\tau}^{t}=0$ при $L\le\tau\le T-1$). Разделив максимальное богатство на $R$ пенсионных периодов, домохозяйство может поддерживать \textbf{то же самое} потребление, что и в трудовой период:
\begin{formulabox}
\begin{align}
s_{t+\tau}^{t} &= -\frac{a_{t+L}^{t}}{R} = -\frac{L}{L+R}\,y^{t}, \tag{6.8}\\
a_{t+\tau}^{t} &= a_{t+L}^{t} + (\tau-L)s_{t+\tau}^{t} = \frac{L(R+L-\tau)}{L+R}\,y^{t}, \qquad L\le\tau\le T-1 \tag{6.9}
\end{align}
\end{formulabox}

\begin{intuition}
Профиль активов на протяжении жизни имеет характерную "треугольную" форму: активы растут линейно на протяжении трудового периода, достигают пика в момент выхода на пенсию (в примере — в 8 раз выше годового дохода при $L=40, R=10$), а затем линейно снижаются до нуля к концу жизни ($\tau=T$). Это графически иллюстрируется на диаграмме, где доход $y_\tau^t\equiv 1$ (горизонтальная линия на уровне 1 до $L=40$, затем 0), потребление $c_\tau^t$ постоянно на уровне 0{,}8 на протяжении всей жизни, сбережения $s_\tau^t$ равны $+0{,}2$ в трудовой период и $-0{,}8$ в пенсионный период, а активы $a_\tau^t$ образуют "треугольник" с пиком в точке $\tau=L=40$ на уровне 8 и спускаются к нулю в точке $\tau=T=50$.
\end{intuition}

\subsection{Агрегирование: от индивидуального к макроуровню}

Предположим, что население растёт непрерывно с темпом $n$, а производительность, определяющая трудовой доход, растёт с темпом $g_A$. Тогда в момент $t$ в экономику входит $n(t)=n_0 e^{nt}$ домохозяйств, и они будут иметь трудовой доход $y(t)=y_0 e^{g_A t}$ на протяжении своего трудового периода.

\begin{intuition}
Важно подчеркнуть: здесь предполагается рост производительности \textbf{от поколения к поколению} (то есть каждая новая когорта входит в экономику с более высоким начальным доходом), а не рост дохода \textbf{в течение} трудового периода одного и того же домохозяйства (напомним, что в рамках одной когорты доход постоянен, см. (6.0)).
\end{intuition}

Агрегированный доход в момент времени $t$, $Y(t)$, — это сумма трудовых доходов всех, кто работает в момент $t$, то есть сумма доходов тех, кто вошёл в экономику от $t-L$ периодов назад до настоящего момента:
\begin{formulabox}
\begin{equation}
Y(t) = \int_{t-L}^{t} y(\tau)n(\tau)\,d\tau = \int_{t-L}^{t} y_0 e^{g_A \tau} n_0 e^{n\tau}\,d\tau \tag{6.10}
\end{equation}
\end{formulabox}

\subsection{Агрегированное потребление}

Агрегированное потребление определяется как сумма расходов на потребление всех живущих поколений (начиная с $t-T$):
\begin{formulabox}
\begin{align}
C(t) &= \int_{t-T}^{t} c(\tau)n(\tau)\,d\tau = \int_{t-T}^{t} \frac{L}{L+R}\,y(\tau)n(\tau)\,d\tau \notag\\
&= \int_{t-L-R}^{t} \frac{L}{L+R}\, y_0 e^{g_A \tau} n_0 e^{n\tau}\,d\tau \tag{6.11}
\end{align}
\end{formulabox}

Разбивая интеграл на две части — для тех, кто ещё работает, и для тех, кто уже вышел на пенсию:
\begin{formulabox}
\begin{align}
C(t) &= \int_{t-L-R}^{t-L} \frac{L}{L+R}\,y(\tau)n(\tau)\,d\tau + \int_{t-L}^{t} \frac{L}{L+R}\,y(\tau)n(\tau)\,d\tau \notag\\
&= \underbrace{\int_{t-L-R}^{t-L} \frac{L}{L+R}\,y(\tau)n(\tau)\,d\tau}_{\text{потребление пенсионеров}} + \underbrace{\frac{L}{L+R}Y(t)}_{\text{потребление работающих}} \tag{6.12}
\end{align}
\end{formulabox}

\begin{intuition}
В то время как работающие домохозяйства потребляют из трудового дохода, пенсионеры потребляют из накопленного богатства. Таким образом, построены \textbf{микрооснования для кейнсианской функции потребления}: агрегированное потребление является линейной функцией агрегированного дохода, а автономное потребление отражает потребление из богатства.
\end{intuition}

\begin{formulabox}
\begin{equation}
C(t) = \alpha Y(t) + \beta A(t) \tag{6.13}
\end{equation}
\end{formulabox}

\subsection{Агрегированные сбережения}

Определим агрегированные сбережения:
\begin{formulabox}
\begin{align}
S(t) &= Y(t) - C(t) \notag\\
&= \underbrace{\int_{t-L}^{t} \frac{R}{L+R}\, y_0 e^{g_A \tau} n_0 e^{n\tau}\,d\tau}_{\text{положительные сбережения работающих}} + \underbrace{\left[-\int_{t-L-R}^{t-L} \frac{L}{L+R}\, y_0 e^{g_A \tau} n_0 e^{n\tau}\,d\tau\right]}_{\text{отрицательные сбережения пенсионеров}} \tag{6.14}
\end{align}
\end{formulabox}

\subsection{Агрегированная средняя склонность к потреблению и сбережению}

Вычисляя интегралы в (6.12) и (6.14), приходим к агрегированной средней склонности к потреблению:
\begin{formulabox}
\begin{equation}
c = \frac{C(t)}{Y(t)} = \frac{L}{L+R}\cdot \frac{1 - e^{-(g_A+n)(L+R)}}{1 - e^{-(g_A+n)L}} \tag{6.15}
\end{equation}
\end{formulabox}

Агрегированная склонность к потреблению \textbf{отличается} от индивидуальной и зависит от темпа роста $g_Y = g_A + n$. Аналогично определим агрегированную (среднюю) склонность к сбережению $s = 1-c$:
\begin{formulabox}
\begin{equation}
s = 1 - \frac{L}{L+R}\cdot \frac{1 - e^{-(g_A+n)(L+R)}}{1 - e^{-(g_A+n)L}} \tag{6.16}
\end{equation}
\end{formulabox}

\begin{intuition}
Агрегированные склонности к потреблению и сбережению зависят от темпа роста экономики. Более низкий темп роста населения может быть компенсирован более быстрым ростом производительности — значение имеет только их сумма $g_Y = g_A+n$.
\end{intuition}

\subsubsection{Агрегированная норма сбережений как функция роста}

Агрегированная средняя склонность к сбережению является \textbf{возрастающей и вогнутой} функцией темпа роста экономики:
$$g_Y = 0\%,\ s=0\%; \qquad g_Y=1\%,\ s=4{,}5\%; \qquad g_Y=3\%,\ s=11{,}1\%; \qquad g_Y=5\%,\ s=15{,}1\%.$$

\begin{intuition}
Более быстрый рост экономики приводит к более высокой агрегированной норме сбережений при той же индивидуальной норме сбережений (в примере — 0{,}2). Если экономика не растёт ($g_Y=0$), агрегированная норма сбережений равна нулю: отрицательные сбережения пенсионеров в точности компенсируют положительные сбережения работающих (все когорты имеют одинаковый размер и одинаковый доход, поэтому сбережения одной когорты зеркально отражаются "растратами" другой). Агрегированные сбережения положительны, если поток положительных сбережений работающих домохозяйств превышает отрицательные сбережения пенсионеров. Это следует:
\begin{itemize}[leftmargin=1.3em]
\item из роста населения (\textbf{эффект Нейссера}, Neisser effect) — работающих когорт становится больше, чем пенсионных;
\item и из роста производительности (\textbf{эффект Бенцеля}, Bentzel effect) — более молодые (и потому более богатые по абсолютному доходу) когорты сберегают больше в абсолютном выражении, чем "проедают" более старые, менее богатые когорты.
\end{itemize}
\end{intuition}

\subsection{Эмпирические свидетельства}

\subsubsection{Modigliani (1986, 1990)}

Оценивая межстрановую зависимость нормы сбережений от темпа роста подушевого располагаемого дохода, Модильяни получил:
$$s = \underset{(4.2)}{0{,}06} + \underset{(6.1)}{1{,}81}\,g_Y, \qquad \bar R^2 = 0{,}37,\quad S.E.=0{,}041.$$
(в скобках под коэффициентами — $t$-статистики). Диаграмма рассеяния по странам ОЭСР (Великобритания, США, Канада, Франция, Германия, Италия, Япония) показывает выраженную положительную связь между темпом роста и нормой сбережений, в целом соответствующую линии "нормы ОЭСР" (O.E.C.D.\ Norm).

\subsubsection{Paxson (1996)}

Диаграмма рассеяния «инвестиции/ВВП» против «роста подушевого ВВП» (десятилетние средние, панель стран, включая США, Великобританию, Тайвань, Таиланд) из работы Paxson C. (1996) ``Savings and Growth: Evidence from Micro Data'', \textit{European Economic Review}, 40(2), pp.\,255--288, также иллюстрирует положительную связь между темпами роста и нормой инвестиций/сбережений на межстрановых данных.

\subsubsection{Экономический рост и сбережения: направление причинности}

Гипотеза жизненного цикла утверждает, что причинность должна идти \textbf{от экономического роста к норме сбережений}, а не наоборот. Эмпирические исследования подтверждают эту гипотезу — см., например, Carroll C.\,D., Weil D.\,N. (1994) ``Saving and Growth: A Reinterpretation''. \textit{Carnegie-Rochester Series on Public Policy}, 40(June), pp.\,133--192. Альтернативное объяснение основано на формировании привычек потребления (habit formation).

\subsection{Возвращение на индивидуальный уровень: критика упрощений}

Модель LCH весьма схематична. Соответствует ли она реальности? Устойчива ли она к изменению (упрощающих) предпосылок?
\begin{itemize}[leftmargin=1.3em]
\item Предполагается нулевая процентная ставка (и нулевая субъективная норма дисконтирования);
\item предполагается, что трудовой доход постоянен на протяжении трудового периода;
\item длительность трудового и пенсионного периодов задаётся экзогенно (продолжительность жизни детерминирована);
\item не учитывается размер домохозяйства (семьи);
\item не вводится пенсионная система.
\end{itemize}

\subsection{Эмпирические профили жизненного цикла (Jappelli--Modigliani, 1998)}

Работа Jappelli T., Modigliani F. (1998) ``The Age-Saving Profile and the Life-Cycle Hypothesis'' даёт эмпирические возрастные профили ключевых переменных по данным домохозяйств, сгруппированных по когортам (по оси абсцисс — возраст от 20 до 80 лет):

\begin{itemize}[leftmargin=1.3em]
\item \textbf{Потребление} (Consumption, Figure 1) растёт от уровня около 20--22 в возрасте 20 лет до пика около 30 в возрасте 45--50 лет, затем плавно снижается до уровня около 16--17 к 80 годам.
\item \textbf{Размер семьи} (Family size, Figure 2) имеет выраженную колоколообразную форму: растёт от чуть более 2 человек в 20 лет до пика почти 4 человек к 45--50 годам, затем снижается до менее 2 человек к 80 годам — это отражает типичный демографический цикл семьи (рождение и взросление детей).
\item \textbf{Трудовой доход и потребление} (Earned Income and Consumption, Figure 3): трудовой доход растёт от около 22 в 20 лет до пика около 31 в 45--50 лет, затем резко падает практически до нуля к 65--70 годам (выход на пенсию), тогда как потребление (сглаженная кривая) остаётся заметно выше трудового дохода в пожилом возрасте — за счёт использования накопленного богатства.
\item \textbf{Совокупные сбережения} (Total saving, Figure 4) положительны и растут до пика около 20--23 в возрасте 35--40 лет, затем снижаются, становятся отрицательными примерно к 65 годам и достигают около $-10$ к 80 годам — классический "треугольный" профиль сбережений LCH.
\item \textbf{Располагаемый доход и потребление} (Disposable Income and Consumption, Figure 5): аналогично Figure 3, но с учётом располагаемого (а не только трудового) дохода — потребление устойчиво превышает располагаемый доход в старших возрастах.
\item \textbf{Частные сбережения} (Private saving, Figure 6) остаются положительными и относительно стабильными (около 5--15) на протяжении почти всей жизни, без выраженного ухода в отрицательную область в старости — в отличие от совокупных сбережений.
\item \textbf{Обязательные сбережения} (Mandatory saving, Figure 7 — пенсионные взносы) демонстрируют чёткий "треугольный" профиль: рост до пика около 14--15 к 35--40 годам, затем снижение и переход в отрицательную область (пенсионные выплаты) до уровня около $-15$ к 80 годам.
\item \textbf{Совокупные сбережения при ставке взносов 20\%} (Total saving, contr.=20\%, Figure 8) сочетает частные и обязательные (пенсионные) сбережения, давая профиль, качественно похожий на профиль Total saving (Figure 4).
\item \textbf{Совокупное богатство} (Total wealth, Figure 9) растёт от около 100--120 в 20 лет до пика около 380--400 к 55--60 годам, затем снижается до около 220--230 к 80 годам — классическая "холмообразная" траектория богатства жизненного цикла.
\item \textbf{Пенсионное богатство} (Pension wealth, Figure 10) имеет аналогичный колоколообразный профиль с пиком около 190--200 к 55--60 годам.
\item \textbf{Частное богатство} (Private wealth, Figure 11) также имеет колоколообразный профиль с пиком около 210--220 к 55--60 годам.
\item \textbf{Изменение богатства} (Change in private wealth и Change in total wealth, Figure 12) положительно и достигает пика (около 10--14) в возрасте 35--45 лет, затем снижается, становится отрицательным примерно после 60 лет и достигает около $-10$ к 80 годам.
\end{itemize}

\begin{intuition}
Эмпирические профили Jappelli--Modigliani в целом \textbf{качественно подтверждают} ключевое предсказание LCH — "холмообразный" (треугольный) профиль накопления и последующего истощения богатства на протяжении жизни, с пиком примерно в предпенсионном возрасте. Вместе с тем эмпирические кривые более гладкие, чем предсказывает упрощённая модель с постоянным доходом (трудовой доход в реальности растёт и падает плавно, а не скачком в момент выхода на пенсию), а размер семьи (Figure 2) явно демонстрирует, что упрощающее предположение о неизменном составе домохозяйства — не более чем аппроксимация.
\end{intuition}

\newpage
\section{Лекция 7. Гипотеза перманентного дохода и buffer-stock saving}

\subsection{Мотивация}

К настоящему моменту рассмотрено уже несколько важных объяснений сбережений:
\begin{itemize}[leftmargin=1.3em]
\item сбережения нужны для накопления производительного капитала;
\item сбережения и заимствования (или потребление из богатства) означают перераспределение неравномерного потока дохода во времени;
\item добровольные или обязательные сбережения на пенсию — особенность жизненного цикла человека (и системы социального обеспечения).
\end{itemize}

Но этот список всё ещё неполон:
\begin{itemize}[leftmargin=1.3em]
\item необходимо глубже изучить сбережения как перераспределение богатства во времени;
\item до сих пор не обсуждались следствия \textbf{неопределённости}, с которой сталкиваются люди.
\end{itemize}

Чтобы выявить другие мотивы сбережений, абстрагируемся от предыдущих концепций накопления капитала и пенсионного богатства.

\subsection{Гипотеза перманентного дохода (PIH): детерминированный случай}

Рассмотрим репрезентативного агента с бесконечным горизонтом жизни. Возьмём уравнение Эйлера и межвременное бюджетное ограничение из Лекции 2:
\begin{formulabox}
\begin{align}
\frac{u'(C_t)}{u'(C_{t+1})} &= \frac{1+r}{1+\rho} \tag{7.1}\\
\sum_{\tau=0}^{\infty} \frac{C_{t+\tau}}{(1+r)^\tau} &= A_t + \sum_{\tau=0}^{\infty} \frac{Y_{t+\tau}^{d}}{(1+r)^\tau} \tag{7.2}
\end{align}
\end{formulabox}

Для упрощения алгебры предположим $r=\rho$. Тогда, независимо от динамики располагаемого дохода, оптимальный выбор — \textbf{постоянное потребление}, $C_{t+1}=C_t=C$.

Подставляя $C_{t+\tau}=C\ \forall\tau$ в (7.2):
$$C\sum_{\tau=0}^{\infty}\frac{1}{(1+r)^\tau} = A_t + \sum_{\tau=0}^{\infty}\frac{Y_{t+\tau}^d}{(1+r)^\tau}.$$

Поскольку $\sum_{\tau=0}^{\infty}(1+r)^{-\tau} = (1+r)/r$, получаем:
\begin{formulabox}
\begin{equation}
C = \frac{r}{1+r}\left(A_t + \sum_{\tau=0}^{\infty}\frac{Y_{t+\tau}^d}{(1+r)^\tau}\right) \overset{\text{def}}{=} Y_t^{P} \tag{7.3}
\end{equation}
\end{formulabox}

\textbf{Перманентный доход} — это аннуитетная (рентная) стоимость совокупного богатства, которое является суммой:
\begin{itemize}[leftmargin=1.3em]
\item финансовых активов $A_t$;
\item и человеческого капитала (дисконтированной суммы будущего потока располагаемого дохода).
\end{itemize}

\begin{intuition}
Эта межвременная логика отличается от кейнсианских предпосылок: потребление определяется \textbf{перманентным}, а не \textbf{текущим} доходом. Человек "сглаживает" не только неравномерность дохода внутри жизненного цикла (как в LCH), но и любую временную/случайную волатильность дохода вообще, ориентируясь на устойчивую, аннуитетную оценку своего совокупного богатства.
\end{intuition}

\subsubsection{Напоминание: аннуитет}

Предположим, что вы покупаете актив (аннуитет), который будет обеспечивать вам постоянную купонную (аннуитетную) выплату $x$ каждый период. Соотношение между ценой аннуитета $X$ и выплатой $x$:
$$X = \frac{x}{1+r} + \frac{x}{(1+r)^2} + \dots = \frac{x}{r} \qquad \text{или} \qquad x = \frac{r}{1+r}X.$$

Такое же соотношение можно записать и для конечного горизонта — оно будет лишь немного сложнее (это полезно каждому: подумайте о банковском кредите и ежемесячном/ежегодном платеже по нему). Определение перманентного дохода в точности воспроизводит эту логику!

\subsubsection{Иллюстрация: конечный горизонт при $r=\rho=0$}

Если расчёты при бесконечном горизонте кажутся странными, рассмотрим простой пример. Пусть потребитель живёт $T$ периодов, имеет некоторые начальные активы $A_0$, и для дальнейшего упрощения $r=\rho=0$:
$$\sum_{\tau=0}^{T-1} u(C_\tau) \to \max_{C_\tau}, \qquad \sum_{\tau=0}^{T-1} C_\tau = A_0 + \sum_{\tau=0}^{T-1} Y_\tau^{d}.$$

Уравнение Эйлера при $r=\rho=0$ даёт $u'(C_t)=u'(C_{t+1})\ \forall t\ \Rightarrow\ C_0=C_1=\dots=C_{T-1}=C$. Тогда:
$$\sum_{\tau=0}^{T-1} C = T\cdot C = A_0 + \sum_{\tau=0}^{T-1}Y_\tau^{d} \quad\Rightarrow\quad C = \frac{1}{T}\left(A_0 + \sum_{\tau=0}^{T-1}Y_\tau^{d}\right).$$

Для произвольного момента $t$ внутри горизонта:
$$C = \frac{1}{T-t}\left(A_t + \sum_{\tau=0}^{T-t-1}Y_{t+\tau}^{d}\right).$$

Сравним этот результат с уравнением (7.3) при $t=0$:
$$C = \frac{r}{1+r}\left(A_0 + \sum_{\tau=0}^{\infty}\frac{Y_\tau^{d}}{(1+r)^\tau}\right) \qquad \longleftrightarrow \qquad C = \frac{1}{T}\left(A_0+\sum_{\tau=0}^{T-1}Y_\tau^{d}\right).$$

\begin{intuition}
Оба выражения говорят об одном и том же: для $T$-периодного горизонта и нулевой процентной ставки \textbf{аннуитетный коэффициент} равен $1/T$ (обратная величина числа оставшихся периодов), а для бесконечного горизонта — $r/(1+r)$. Но следует быть осторожным при интерпретации перманентного дохода как \textbf{среднего} дохода — это лишь частный случай (справедливый при $r=\rho=0$), а не общее правило.
\end{intuition}

\subsection{Транзиторный доход}

Определим разницу между перманентным и текущим доходом как \textbf{транзиторный доход}:
\begin{formulabox}
\begin{equation}
Y_t^{T} \overset{\text{def}}{=} Y_t^{d} - Y^{P} \tag{7.4}
\end{equation}
\end{formulabox}
где
$$Y^{P} = \frac{r}{1+r}\left(A_t + \sum_{\tau=0}^{\infty}\frac{Y_{t+\tau}^{d}}{(1+r)^\tau}\right).$$

Поскольку $Y^{P}=C$, транзиторный доход — это, по сути, то же самое, что и \textbf{сбережения}: он позволяет перераспределять ресурсы во времени — сберегать, когда текущий доход превышает перманентный, и занимать (или потреблять из богатства), когда текущий доход ниже перманентного.

\subsection{Проверьте себя: пример с выигрышем в лотерею}

Предположим, вы выиграли \$1\,000\,000. Оценим оставшийся срок жизни в 50 лет. Следуя PIH, вам следует увеличить \textbf{годовые} расходы всего на \$20\,000 (\$1\,000\,000/50).

\begin{intuition}
Естественная реакция многих людей — "НЕТ!!! Я хочу потратить миллион СЕЙЧАС!!!" Возникает вопрос: действительно ли вы ведёте себя как "умный рациональный" агент из модели PIH? Этот пример иллюстрирует расхождение между теоретическим предсказанием строгой версии PIH и интуитивным поведением реальных людей — тема, к которой мы вернёмся при обсуждении ограничений ликвидности и предосторожных сбережений.
\end{intuition}

\subsection{Проверьте других: соответствуют ли "поведенческие паттерны" PIH?}

\begin{itemize}[leftmargin=1.3em]
\item Бедные люди сберегают меньшую долю своего дохода, чем богатые, поскольку им приходится удовлетворять минимальный уровень жизни, покупать предметы первой необходимости и т.д.
\item Некоторые люди сберегают слишком мало, чтобы "не отставать от соседей" (keep up with the Joneses).
\end{itemize}

\subsection{PIH и оценённая функция потребления}

\subsubsection{Кейнсианский крест и панельные данные}

Кейнс предполагал, что функция потребления стабильна, и что более высокий доход ведёт к более высокой норме сбережений. Эконометрические оценки на панельных данных домохозяйств подтверждают эту гипотезу — оценённая линия потребления по индивидуальным (кросс-секционным) данным имеет положительный наклон меньше единицы и положительный свободный член.

\subsubsection{Загадка Кузнеца (Kuznets puzzle)}

Однако агрегированные временны́е ряды показывают \textbf{иное} соотношение: в отличие от кросс-секционных данных, долгосрочные временны́е ряды потребления и дохода демонстрируют практически пропорциональную зависимость (соотношение $C/Y$ остаётся стабильным во времени), в то время как кросс-секционные данные предполагают убывающую среднюю склонность к потреблению по мере роста дохода. Это расхождение известно как \textbf{загадка Кузнеца}.

Кроме того, различия в оценках функций потребления для разных групп домохозяйств с трудом объясняются в рамках кейнсианской традиции.

\subsubsection{Формальное объяснение через PIH}

Пусть оценивается регрессия:
\begin{formulabox}
\begin{equation}
C_i = a + bY_i + e_i \tag{7.$*$}
\end{equation}
\end{formulabox}

Оценка наклона:
$$\hat b = \frac{\mathrm{Cov}(Y,C)}{\mathrm{Var}(Y)} = \frac{\mathrm{Cov}(Y^{P}+Y^{T},\,Y^{P})}{\mathrm{Var}(Y^{P}+Y^{T})} = \frac{\mathrm{Var}(Y^{P})}{\mathrm{Var}(Y^{P})+\mathrm{Var}(Y^{T})},$$
$$\hat a = \bar C - \hat b\bar Y = \bar Y^{P} - \hat b(\bar Y^{P}+\bar Y^{T}) = (1-\hat b)\bar Y^{P}.$$

Наклон оценённой функции потребления зависит от \textbf{относительной вариации} $Y^{P}$ и $Y^{T}$. Рост текущего дохода ведёт к росту потребления только в той мере, в какой он связан с ростом перманентного дохода.

\begin{intuition}
\begin{itemize}[leftmargin=1.3em]
\item Если $\mathrm{Var}(Y^{T}) \gg \mathrm{Var}(Y^{P})$, то изменения текущего дохода в основном отражают изменения транзиторного дохода. Следовательно, они оказывают относительно слабое влияние на потребление ($\hat b<1$, $\hat a>0$).
\item Если $\mathrm{Var}(Y^{T}) \ll \mathrm{Var}(Y^{P})$, то изменения текущего дохода в основном отражают изменения перманентного дохода. Следовательно, потребление меняется "один к одному" вместе с текущим доходом.
\end{itemize}
Для отдельного домохозяйства бо́льшая часть вариации дохода обусловлена личными обстоятельствами (возраст, занятость, карьера и т.д.), то есть $\mathrm{Var}(Y^{T})\gg\mathrm{Var}(Y^{P})$. Поэтому оценки на панельных данных дают $\hat b<1$, $\hat a>0$. На агрегированном уровне (в динамике) бо́льшая часть изменений дохода отражает экономический рост (устойчивую компоненту), поэтому оценки по временны́м рядам дают $\hat b\approx 1$, $\hat a\approx 0$. Известный эмпирический факт: "чёрные" и "белые" домохозяйства (в терминологии оригинальных американских исследований) имеют схожие предельные склонности к потреблению, но разные уровни перманентного дохода $Y^{P}$ — что объясняет кажущиеся различия в оценённых функциях потребления между группами.
\end{intuition}

\subsection{PIH при неопределённости будущих доходов}

В реальности домохозяйства не знают своего будущего дохода наверняка. Они максимизируют \textbf{ожидаемую} пожизненную полезность при бюджетном ограничении в терминах ожидаемой приведённой стоимости:
\begin{formulabox}
\begin{align}
E_0\left[\sum_{t=0}^{\infty} \frac{u(C_t)}{(1+\rho)^t}\right] &\to \max_{C_t} \tag{7.5}\\
E_t\left[\sum_{\tau=0}^{\infty}\frac{C_{t+\tau}}{(1+r)^\tau}\right] &= A_t + E_t\left[\sum_{\tau=0}^{\infty}\frac{Y_{t+\tau}^{d}}{(1+r)^\tau}\right] \tag{7.6}
\end{align}
\end{formulabox}

В этом случае уравнение Эйлера принимает форму:
\begin{formulabox}
\begin{equation}
\frac{u'(C_t)}{E_t[u'(C_{t+1})]} = \frac{1+r}{1+\rho} \tag{7.7}
\end{equation}
\end{formulabox}

Вновь предположим $r=\rho$. Тогда:
\begin{formulabox}
\begin{equation}
u'(C_t) = E_t[u'(C_{t+1})] \tag{7.8}
\end{equation}
\end{formulabox}

\subsection{Гипотеза случайного блуждания (Random Walk Hypothesis)}

Предположим, что $E_t[u'(C_{t+1})] = u'(E_t[C_{t+1}])$. Это соответствует квадратичной функции полезности (линейной предельной полезности). Это предположение известно как \textbf{эквивалентность определённости} (certainty equivalence): агент выбирает потребление так, как будто он знает истинные значения будущих доходов.

Это упрощает уравнение Эйлера (7.8):
$$u'(C_t) = E_t[u'(C_{t+1})] = u'(E_t[C_{t+1}]) \quad \Rightarrow$$
\begin{formulabox}
\begin{equation}
E_t[C_{t+1}] = C_t \qquad \text{или} \qquad C_{t+1} = C_t + \varepsilon_{t+1},\quad \varepsilon_{t+1}\sim i.i.d.(0,\sigma^2) \tag{7.9}
\end{equation}
\end{formulabox}

Это и есть так называемая \textbf{гипотеза случайного блуждания}: изменение потребления в каждом периоде — случайная величина с нулевым математическим ожиданием. Отсюда следует, что \textbf{наилучший прогноз} любого будущего потребления — это \textbf{текущее} потребление: $E_t[C_{t+\tau}] = C_t\ \forall\tau$.

\begin{intuition}
Гипотеза случайного блуждания Холла (Hall, 1978) — одно из самых знаменитых и проверяемых следствий теории рациональных ожиданий в макроэкономике. Она означает, что \textbf{никакая} доступная в момент $t$ информация (включая прошлые значения дохода, потребления и т.д.), кроме самого $C_t$, не должна иметь предсказательной силы для будущего изменения потребления $C_{t+1}-C_t$ — иначе это означало бы, что агент не полностью использует всю доступную ему информацию при формировании ожиданий, что противоречит гипотезе рациональных ожиданий.
\end{intuition}

\subsubsection{Определение перманентного дохода при неопределённости}

Аналогично (7.3), приходим к похожему определению перманентного дохода, но теперь — как условного математического ожидания:
\begin{formulabox}
\begin{equation}
C_t = \frac{r}{1+r}\left(A_t + E_t\left[\sum_{\tau=0}^{\infty}\frac{Y_{t+\tau}^{d}}{(1+r)^\tau}\right]\right) \overset{\text{def}}{=} Y_t^{P} \tag{7.9$'$}
\end{equation}
\end{formulabox}

Но теперь перманентный доход (а значит, и потребление) \textbf{не постоянен} — он следует некоторому стохастическому процессу. Что определяет случайное изменение потребления?

\subsubsection{Вывод изменения потребления}

Запишем (7.9$'$) для $t$ и для $t+1$:
$$C_t = \frac{r}{1+r}\left(A_t + E_t\left[\sum_{\tau=0}^{\infty}\frac{Y_{t+\tau}^{d}}{(1+r)^\tau}\right]\right), \qquad C_{t+1} = \frac{r}{1+r}\left(A_{t+1} + E_{t+1}\left[\sum_{\tau=0}^{\infty}\frac{Y_{t+1+\tau}^{d}}{(1+r)^\tau}\right]\right).$$

Используя динамическое бюджетное ограничение, можно показать, что:
\begin{formulabox}
\begin{equation}
C_{t+1} - C_t = \Delta C_{t+1} = \frac{r}{1+r}\sum_{\tau=0}^{\infty}\frac{E_{t+1}[Y_{t+1+\tau}^{d}] - E_t[Y_{t+1+\tau}^{d}]}{(1+r)^\tau} \tag{7.10}
\end{equation}
\end{formulabox}

Изменение потребления между $t$ и $t+1$ определяется \textbf{изменением ожиданий} относительно будущих доходов. Это изменение случайно постольку, поскольку случайно изменение самих ожиданий.

\begin{intuition}
\textbf{Только новая информация} о будущих доходах ведёт к изменению перманентного дохода (потребления). Если домохозяйство не получает никаких новостей, потребление остаётся неизменным. Из этого следует, в частности, что только \textbf{неожиданные} изменения налогов могут повлиять на потребление — потребление меняется в момент объявления новой налоговой политики, а не в момент её фактического вступления в силу (если она была анонсирована заранее и уже учтена в ожиданиях). Временно́е (transitory) изменение налоговой политики должно оказывать на потребление \textbf{меньшее} влияние, чем постоянное (permanent) изменение — поскольку оно меняет приведённую стоимость будущих располагаемых доходов в меньшей степени. Если прошлые значения дохода позволяют прогнозировать будущие доходы (то есть доход не является чистым случайным блужданием), то изменение потребления \textbf{не является} чисто случайным — это явление называется \textbf{избыточной чувствительностью} (excess sensitivity) потребления к предсказуемым изменениям дохода, и оно рассматривается как эмпирическое отклонение от строгой RWH.
\end{intuition}

\subsection{Предусмотрительный потребитель (Prudent Consumer)}

Хотя эквивалентность определённости — удобная концепция выбора в условиях неопределённости, она \textbf{игнорирует влияние изменения рисков} — что представляется нереалистичным.

Рассмотрим пример: пусть нужно выбрать уровень потребления сегодня в условиях неопределённости дохода в следующем периоде. Рассмотрим два сценария:
\begin{itemize}[leftmargin=1.3em]
\item \textbf{Сценарий \#1}: с вероятностью 0{,}5 доход составит \$1200, и только \$800 в противном случае;
\item \textbf{Сценарий \#2}: с вероятностью 0{,}5 доход составит \$1500, и только \$500 в противном случае.
\end{itemize}

При эквивалентности определённости эти сценарии \textbf{идентичны}, поскольку оба предполагают одинаковый ожидаемый доход \$1000. Но выбрали бы вы одинаковый уровень потребления в обоих сценариях, или проявили бы предусмотрительность (prudence)? Как смоделировать эту предусмотрительность?

\subsubsection{Моделирование мотива предусмотрительности}

Чтобы ввести мотив предусмотрительности, необходимо предположить $u'''(\cdot) > 0$. В этом случае неравенство Йенсена даёт:
\begin{formulabox}
\begin{align}
E_t[u'(C_{t+1})] &> u'(E_t[C_{t+1}]) \notag\\
E_t[u'(C_{t+1})]\big|_{\sigma_2^2} &> E_t[u'(C_{t+1})]\big|_{\sigma_1^2}, \qquad \sigma_2^2 > \sigma_1^2 \tag{7.$**$}
\end{align}
\end{formulabox}

То есть \textbf{ожидаемая} предельная полезность выше, чем предельная полезность от \textbf{ожидаемого} потребления, причём чем выше риск, тем выше ожидаемая предельная полезность.

\begin{intuition}
Поскольку предельная полезность является убывающей функцией потребления, это означает, что более высокий риск подразумевает более низкое $C_{t+1}$ (при заданном уравнении Эйлера, требующем равенства $u'(C_t)$ и ожидаемой $E_t[u'(C_{t+1})]$). Иными словами, сталкиваясь с более высоким риском в отношении завтрашнего дня, потребитель сберегает больше сегодня. Графически: функция $u'(C)$ выпукла (поскольку $u'''>0$), поэтому при более широком разбросе возможных значений $C_{t+1}$ вокруг $E_t[C_{t+1}]$ среднее значение $u'(C_{t+1})$ (то есть $E_t[u'(C_{t+1})])$, вычисленное по выпуклой функции, оказывается выше, чем значение функции в средней точке $u'(E_tC_{t+1})$, причём это превышение тем больше, чем больше дисперсия $\sigma^2$.
\end{intuition}

\subsection{Предосторожные сбережения (Precautionary Saving)}

Мотив предусмотрительности вводит ещё одну причину для сбережений — так называемые \textbf{предосторожные сбережения} (precautionary saving). Это отличается от ранее обсуждавшихся мотивов (сбережения для накопления капитала и увеличения будущего производства, сбережения для перераспределения ресурсов во времени, сбережения на пенсионный период). Но это не означает, что RWH/PIH, основанные на эквивалентности определённости, полностью лишены смысла.

Предусмотрительность порождает относительно небольшой дополнительный объём удерживаемых активов, так называемые \textbf{буферные сбережения} (buffer-stock saving).

\begin{intuition}
Этот мотив ослабевает по мере накопления активов (чем больше у домохозяйства уже накоплено "подушки безопасности", тем меньше предельная выгода от дополнительной предосторожности). Это подразумевает некоторый \textbf{восходящий тренд} в потреблении: молодое домохозяйство без достаточного богатства должно потреблять меньше (сберегая "про запас"), но по мере накопления богатства оно будет потреблять больше в будущем. Этот эффект действует \textbf{в противоположном направлении} по отношению к нетерпеливости (высокому $\rho$), которая, наоборот, порождает нисходящий тренд потребления.
\end{intuition}

\subsection{Ограничения ликвидности}

На протяжении всего анализа предполагалось, что домохозяйства могут сберегать и занимать по одной и той же процентной ставке — единственным ограничением на заимствование (или чрезмерное накопление) было условие NPG.

В реальности люди часто сталкиваются с \textbf{ограничениями ликвидности}:
\begin{itemize}[leftmargin=1.3em]
\item процент по сбережениям ниже, чем ставка по кредиту, с которой они сталкиваются;
\item некоторые люди вообще не имеют доступа к кредиту, что означает $A_t \ge 0$.
\end{itemize}

\begin{intuition}
Столкновение с ограничениями ликвидности (сейчас или в будущем) подразумевает \textbf{дополнительный мотив} для буферных сбережений: если в периоды низкого дохода трудно занять средства, лучше заранее накопить дополнительное богатство, чтобы иметь возможность сглаживать потребление в будущем.

Буферные сбережения — это \textbf{корень избыточных сбережений во время кризиса}: в условиях экономической неопределённости (рост риска потери работы, ужесточение условий кредитования) домохозяйства резко наращивают предосторожные и буферные сбережения, что ведёт к сокращению совокупного спроса. Это и есть теоретическое обоснование призыва к \textbf{кейнсианской политике} стимулирования спроса в кризис (вспомним аргументы Кругмана, обсуждавшиеся в связи с кризисом 2008 г. в Лекциях 1--2).
\end{intuition}

\subsection{Краткое резюме ключевых результатов (Лекции 6--7)}

\begin{itemize}[leftmargin=1.3em]
\item \textbf{Гипотеза жизненного цикла (LCH)}: при постоянном трудовом доходе на протяжении $L$ трудовых периодов и его отсутствии в течение $R$ пенсионных периодов оптимальное (сглаженное) потребление составляет $c^t = \frac{L}{L+R}y^t$; индивидуальный профиль активов имеет "треугольную" форму с пиком в момент выхода на пенсию.
\item При агрегировании по растущей и продуктивно развивающейся экономике LCH даёт \textbf{микрооснования кейнсианской функции потребления} $C(t)=\alpha Y(t)+\beta A(t)$; агрегированная норма сбережений положительна и растёт вместе с темпом роста экономики $g_Y=g_A+n$ (эффекты Нейссера и Бенцеля), что подтверждается межстрановыми эмпирическими данными (Modigliani, Paxson) и возрастными профилями богатства (Jappelli--Modigliani).
\item \textbf{Гипотеза перманентного дохода (PIH)}: потребление определяется \textbf{перманентным доходом} — аннуитетной стоимостью совокупного богатства (финансовых активов плюс человеческого капитала), а не текущим доходом; это объясняет "загадку Кузнеца" — расхождение между кросс-секционными и временны́ми оценками функции потребления.
\item При неопределённости и эквивалентности определённости PIH порождает \textbf{гипотезу случайного блуждания} потребления (Hall, 1978): потребление меняется только в ответ на новую информацию о будущих доходах; отсюда — относительная неэффективность анонсированных заранее (или временных) изменений налоговой политики.
\item \textbf{Мотив предусмотрительности} ($u'''>0$) порождает \textbf{предосторожные (буферные) сбережения}: более высокий риск будущего дохода ведёт к более высоким сбережениям сегодня; этот мотив ослабевает по мере накопления богатства.
\item \textbf{Ограничения ликвидности} усиливают мотив буферных сбережений и лежат в основе резкого роста избыточных сбережений и падения спроса во время кризисов — что обосновывает необходимость кейнсианской антициклической политики.
\end{itemize}
\end{document}